% AxiomForge manuscript — template version 1.0.0.
% Build: latexmk -pdf main.tex   (see README.md)
\documentclass[11pt]{amsart}
\makeatletter
% AxiomForge manuscript preamble — template version 1.0.0.
% This file is copied into every contribution and is owned by it afterwards.
% It fixes the typographic identity of the series; it defines no mathematical notation
% (each paper keeps its own notation in macros.tex).

% ---------------------------------------------------------------- page and fonts
\usepackage[a4paper,hmargin=31mm,top=34mm,bottom=36mm,headsep=10mm,footskip=12mm]{geometry}
\usepackage[T1]{fontenc}
\usepackage[utf8]{inputenc}
\usepackage{lmodern}
\usepackage{amsmath,amssymb,amsthm,mathtools}
\usepackage[final,protrusion=true,expansion=true]{microtype}
\usepackage{booktabs}
\usepackage{array}
\usepackage[shortlabels]{enumitem}
\setlist{itemsep=0.2em,topsep=0.4em}
\usepackage{xcolor}
\linespread{1.04}
\setlength{\emergencystretch}{1.5em}
\allowdisplaybreaks[1]
\numberwithin{equation}{section}
\setcounter{tocdepth}{2}
% Table of contents: sections (and their page numbers) in bold with a little space above
% them, subsections indented in normal weight; after \appendix only the (bold) section
% titles are listed, so the appendices never crowd the contents.
\def\l@section{\@tocline{1}{2.5pt plus 1pt minus 2pt}{0pt}{}{\bfseries\let\@tocpagenum\af@tocpagenumbf}}
\def\af@tocpagenumbf#1{\hss{\bfseries #1}}
\def\l@subsection{\@tocline{2}{0pt}{1.5em}{2.3em}{}}
\g@addto@macro\appendix{\addtocontents{toc}{\protect\setcounter{tocdepth}{1}}}
% Every appendix starts on a new page; the unnumbered back matter (\afbacksection) keeps the
% original \section and flows normally after the last appendix.
% (Local implementation through a hook, so that \section itself is not redefined.)
\newif\ifaf@appendixpage
\AddToHook{cmd/section/before}{\ifaf@appendixpage\clearpage\fi}
\g@addto@macro\appendix{\global\af@appendixpagetrue}

% ---------------------------------------------------------------- links
\definecolor{afblue}{rgb}{0.08,0.18,0.42}
\usepackage[pdfusetitle,bookmarks=true,bookmarksnumbered=true,bookmarksopen=true,
  colorlinks=true,linkcolor=afblue,citecolor=afblue,urlcolor=afblue,
  pdfpagelabels=true,pdfdisplaydoctitle=true]{hyperref}
\usepackage[capitalise,noabbrev]{cleveref}
% Local addition: unique PDF anchors for theorem-like environments numbered within sections.
\hypersetup{hypertexnames=false}

% ---------------------------------------------------------------- theorem environments
% thmtools keeps cleveref names correct for environments sharing one counter.
\usepackage{thmtools}
\declaretheorem[style=plain,numberwithin=section,name=Theorem]{theorem}
\declaretheorem[style=plain,sibling=theorem,name=Proposition]{proposition}
\declaretheorem[style=plain,sibling=theorem,name=Lemma]{lemma}
\declaretheorem[style=plain,sibling=theorem,name=Corollary]{corollary}
\declaretheorem[style=plain,numbered=no,name=Theorem]{theorem*}
\declaretheorem[style=definition,sibling=theorem,name=Definition]{definition}
\declaretheorem[style=definition,sibling=theorem,name=Example]{example}
\declaretheorem[style=remark,sibling=theorem,name=Remark]{remark}
\crefname{equation}{}{}
\Crefname{equation}{Equation}{Equations}

% ---------------------------------------------------------------- series metadata
% \afidentifier{AF-0001}, \afversion{1.0.0}, \afdate{25 November 2026},
% \afrepository{https://github.com/axiomforge-math/mathematics/tree/main/results/...},
% \afdisclosure{short AI-assisted research disclosure},
% \afstatus{verification status} (local addition; optional)
\newcommand{\afidentifier}[1]{\gdef\af@identifier{#1}}
\newcommand{\afversion}[1]{\gdef\af@version{#1}}
\newcommand{\afdate}[1]{\gdef\af@date{#1}}
\newcommand{\afrepository}[1]{\gdef\af@repository{#1}}
\newcommand{\afdisclosure}[1]{\gdef\af@disclosure{#1}}
\newcommand{\afcorrespondence}[1]{\gdef\af@correspondence{#1}}
\gdef\af@correspondence{}
% Local addition: verification status of the contribution, printed as a 'Status' row.
\newcommand{\afstatus}[1]{\gdef\af@status{#1}}
\gdef\af@status{}

% The publication block printed under the abstract on the title page.
\newcommand{\afpublicationblock}{%
  \begingroup\footnotesize\noindent
  \begin{tabular}{@{}l@{\quad}p{0.72\textwidth}@{}}
    \textit{Identifier} & \af@identifier\quad(version \af@version, \af@date)\\[1pt]
    \ifx\af@status\@empty\else\textit{Status} & \af@status\\[1pt]\fi
    \textit{Repository} & \url{\af@repository}\\[1pt]
    \ifx\af@correspondence\@empty\else\textit{Correspondence} & \af@correspondence\\[1pt]\fi
    \textit{Disclosure} & \af@disclosure
  \end{tabular}\par
  \endgroup}

% Front matter: title page, then contents, then the body on a fresh page.
\newcommand{\affrontmatter}{%
  \maketitle
  \vspace*{\fill}
  \afpublicationblock
  \thispagestyle{plain}%
  \clearpage
  {\hypersetup{linkcolor=black}\fontsize{9}{10.6}\selectfont\tableofcontents}
  \clearpage}

% Unnumbered back-matter sections (amsart lists starred sections in the contents).
\newcommand{\afbacksection}[1]{\global\af@appendixpagefalse\section*{#1}}

\makeatother
% Notation local to this paper.
\newcommand{\N}{\mathbb{N}}
\newcommand{\Z}{\mathbb{Z}}
\newcommand{\Q}{\mathbb{Q}}
\newcommand{\R}{\mathbb{R}}

% special functions and the target numbers
\DeclareMathOperator{\Ti}{Ti}
\newcommand{\Li}{\operatorname{Li}}

% general operators
\DeclareMathOperator*{\Res}{Res}
\DeclareMathOperator{\diag}{diag}
\newcommand{\pos}[1]{\left(#1\right)_{+}}
\newcommand{\floor}[1]{\left\lfloor #1\right\rfloor}
\newcommand{\vG}[1]{v_{#1}^{\mathrm{G}}}
\newcommand{\abs}[1]{\left\lvert #1\right\rvert}
\newcommand{\one}{\mathbf{1}}

% objects of the construction
\newcommand{\Rsp}{\mathcal{R}_K}              % rational functions with poles at the x_k
\newcommand{\mom}{\mathcal{S}_w}              % rational moment sum
\newcommand{\nrm}{\mathsf{S}}                 % normalising factor
\newcommand{\Gen}{\mathcal{G}}                % general construction (q >= 83)
\newcommand{\Band}{\mathcal{B}}               % band construction (q = 81)
\newcommand{\Pa}{\mathcal{A}}                 % parity sum A(t)
\newcommand{\rest}[1]{#1\langle z\rangle}     % restricted power series
\newcommand{\aconst}{\mathfrak{a}}            % arithmetic constant
\newcommand{\bconst}{\mathfrak{b}}            % analytic constant
\newcommand{\Iprof}{I}                        % ordinary-prime profile integral
\newcommand{\bvec}{\mathbf{b}}                % binomial column vector
\newcommand{\vvec}{\mathbf{v}}                % monomial column vector
\newcommand{\diff}{\,\mathrm{d}}

% title pieces (robust, so that amsart's upper-casing of the title leaves them unchanged)
\DeclareRobustCommand{\titlechi}{$\chi_2(2/q)$}
\DeclareRobustCommand{\titleti}{$\mathrm{Ti}_2(2/q)$}
\DeclareRobustCommand{\titleq}{$q\ge 81$}


\title[Irrationality of Legendre chi and inverse tangent values]%
{\texorpdfstring{Irrationality of \titlechi{} and \titleti{} for odd \titleq:\\ a candidate proof}%
{Irrationality of chi2(2/q) and Ti2(2/q) for odd q >= 81: a candidate proof}}
\author{AxiomForge}
\urladdr{https://github.com/axiomforge-math/mathematics}
\subjclass[2020]{Primary 11J72; Secondary 33B30}
\keywords{Legendre chi function, inverse tangent integral, irrationality, Hankel determinants,
Pad\'e-type approximation, $p$-adic valuations, logarithmic energy}
\date{26 November 2026}

\afidentifier{AF-0003}
\afversion{0.1.0}
\afdate{26 November 2026}
\afrepository{https://github.com/axiomforge-math/mathematics/tree/main/results/AF-0003-legendre-chi-inverse-tangent-irrationality}
\afcorrespondence{\href{mailto:axiomforge.math@gmail.com}{axiomforge.math@gmail.com}}
\afstatus{Preprint (candidate proof). Not reviewed by an independent human expert; the main theorems are
not formally verified.}
\afdisclosure{AI-assisted research. The candidate arguments were produced by an AI system; their
reconstruction, the audit computations and this manuscript were produced by another AI system
under the direction of AxiomForge. See the AI-assisted research disclosure at the end.}

\begin{document}

\begin{abstract}
Let $\chi_2(x)=\sum_{j\ge0}x^{2j+1}/(2j+1)^2$ be Legendre's chi function and
$\Ti_2(x)=\sum_{j\ge0}(-1)^jx^{2j+1}/(2j+1)^2$ the inverse tangent integral. We present a
candidate proof that $\chi_2(2/q)$ and $\Ti_2(2/q)$ are irrational for every odd integer
$q\ge 81$, and hence also at $-2/q$. The best previously claimed range, in an unrefereed and
primarily AI-generated preprint of Krisman, is odd $q\ge 175$ for $\chi_2(2/q)$. The criterion of
David, Hirata-Kohno and Kawashima applies at $x=-4/q^2$ as well as at $x=4/q^2$ and should give a
threshold of the same order for $\Ti_2(2/q)$; we have not computed it and have not found it
stated. The rational approximations are Hankel determinants of a
moment functional with poles at the negative half-integers, as in Krisman's work, with two
changes: the number of poles exceeds the size of the determinant in the ratio $7/6$, and for
$q=81$ a band of double zeros is placed on the positive axis. The proof needs three estimates
that are new in this setting: a bound at the prime $2$ that separates the interpolation nodes by
parity, a residue-class profile for the ordinary primes valid for any ratio of poles to rows,
and a nonvanishing argument at prime indices that reduces the determinant to three fixed local
Hankel determinants. The finite inequalities on which the argument rests have been recomputed
by second implementations, also written by an AI system, that do not import the code that
produced them. The argument has not been checked
by an independent human expert and its main theorems are not formally verified, so the result
should be read as a candidate.
\end{abstract}

\affrontmatter

\section{Introduction}\label{sec:introduction}

For $\abs{x}\le 1$ let
\[
  \chi_2(x)=\sum_{j\ge0}\frac{x^{2j+1}}{(2j+1)^2}=\tfrac12\bigl(\Li_2(x)-\Li_2(-x)\bigr),
  \qquad
  \Ti_2(x)=\sum_{j\ge0}\frac{(-1)^jx^{2j+1}}{(2j+1)^2}=\operatorname{Im}\Li_2(ix)
\]
be Legendre's chi function and the inverse tangent integral. This paper presents a proof of the
following statement.

\begin{theorem}\label{thm:main}
For every odd integer $q\ge 81$ the numbers $\chi_2(2/q)$ and $\Ti_2(2/q)$ are irrational.
Since $\chi_2$ and $\Ti_2$ are odd, the same holds for $\chi_2(-2/q)$ and $\Ti_2(-2/q)$.
\end{theorem}

\emph{Status.} The proof is a candidate. It was reconstructed and checked by an AI system from
candidate arguments produced by another AI system (see the disclosure at the end); no
independent human expert has reviewed it. In Lean only the final deduction is formalised, and
only conditionally on the three estimates of \cref{thm:three-estimates}. All finite
inequalities used in the proof have been recomputed by second implementations, also written by an
AI system, that do not import the code that produced them (Appendices~\ref{app:certificates}
and~\ref{app:consistency}); this checks the numbers, not the asymptotic arguments, and the two
implementations share inputs and authorship, so it is not an independent verification. \Cref{subsec:step-status} gives the status of each
step, and \cref{rem:scrutiny} lists the points that most need independent scrutiny.

\subsection{Background}\label{subsec:background}

For even $q$ one has $\chi_2(2/q)=\chi_2(1/(q/2))$ and $\Ti_2(2/q)=\Ti_2(1/(q/2))$, so the
interest of \cref{thm:main} lies in odd $q$, where $2/q$ is not the reciprocal of an integer.
We have not read Hata's paper~\cite{Hata1990} ourselves. According to Krisman
\cite[\S1.2]{Krisman2026chi2}, whose reading we have not checked, Hata's linear independence
measures for shifted polylogarithms are stated at points $1/Q$ with $Q\in\Z$, so that for odd $q$
the relevant points $\pm4/q^2$ are not covered. The linear independence criterion of David,
Hirata-Kohno and Kawashima~\cite{DHK2022Acta} applies to the shifted dilogarithm
$\sum_{n\ge1}x^n/(n-\frac12)^2$ at rational points. As reported in \cite[\S1.2]{Krisman2026chi2},
Krisman evaluates it at $x=4/q^2$ and obtains the irrationality of $\chi_2(2/q)$ for
$q\ge 11\,795$ with the criterion as printed, and for $q\ge23\,587$ with a weaker form that he
believes the proof supports; these specialisations are his, not those of the authors of the
criterion. The point $x=-4/q^2$, which gives $\Ti_2(2/q)$, is also a rational point, so the same
criterion should give a threshold of the same order for $\Ti_2(2/q)$; we have not computed it and
have not found it stated. We did not search the classical literature on linear independence of
$1,\Li_2(x),\Li_2(-x)$ (Nikishin, Chudnovsky, Hata, Rhin--Viola, Marcovecchio and others)
ourselves; ``new'' below always means ``not found in the sources we read''.

The closest work is Krisman's preprint~\cite{Krisman2026chi2}, which claims proofs that
$\chi_2(1/q)\notin\Q$ for $q\ge3$, $q\ne4$, that $\Ti_2(1/q)\notin\Q$ for $q\ge5$, and that
$\chi_2(2/q)\notin\Q$ for odd $q\ge175$; it does not treat $\Ti_2(2/q)$ for odd $q$. That
preprint is labelled as primarily AI-generated and has not been refereed. We read version~2 and
checked against it the statement of each numbered result cited in this paper (Lemmas~2.1, 4.2, 5.1, 5.2,
6.4, 7.1, 11.1; Propositions~3.1, 11.2; Theorems~3.2, 4.5, 5.3, 8.1; Remark~6.7; \S1.2, \S12);
this was done by the AI system that revised this manuscript, not by a person, and it checks
statements, not proofs. We use none of its theorems as an input, and every estimate we need is
proved here. Its method, a Hankel-determinant construction
in the spirit of Zudilin's determinantal approach~\cite{Zudilin2017}, is the starting point of
ours: Krisman takes the Hankel determinants of the functional
$g\mapsto\sum_{t\ge1}w^t\,(g(u)\,u/D(u))'|_{u=t}$, where $D(u)=\prod_{k<h}(u+k+\frac12)$ has as
many half-integer poles as the size $h$ of the determinant, at $w=4/q^2$. In his \S12 he
states, numerically and assuming that $\frac{133}{60}$ is the true generic constant, that this
construction does not reach odd $q$ below about $87$: at the prime $2$ its leading coefficient
alone costs $\log2$ per $h^2$ \cite[Prop.~11.2]{Krisman2026chi2}, and with this cost and the
equilibrium constant the margin is positive only for odd $q\lesssim87$. We take this as a
reason to change the determinant rather than to sharpen the estimates for his.

\subsection{The constructions}\label{subsec:constructions}

Fix an odd integer $q\ge5$ (the hypothesis $q\ge5$ is a convenience; only $W<1$ is needed, that
is $q\ge3$, and the main theorem concerns $q\ge81$) and a sign $\sigma\in\{1,-1\}$, and put
\begin{equation}\label{eq:basic-notation}
  y=\frac q2,\qquad w=\frac{\sigma}{y^2}=\frac{4\sigma}{q^2},\qquad W=\abs w,\qquad
  \xi=\sum_{j\ge0}\frac{\sigma^jy^{-2j-1}}{(2j+1)^2},
\end{equation}
so that $\xi=\chi_2(2/q)$ for $\sigma=1$ and $\xi=\Ti_2(2/q)$ for $\sigma=-1$. For integers
$h\ge1$, $K\ge h$ and $A\ge0$ let
\[
  x_k=-k-\tfrac12\ (0\le k<K),\qquad D_K(u)=\prod_{k<K}(u-x_k),\qquad
  \varphi(u)=\frac{u\prod_{j=1}^{A}(u-j)^2}{D_K(u)}.
\]
In \cref{sec:pencils} we define an affine functional $\mu_X$ on the rational functions with
simple poles at the $x_k$ such that $\mu_\xi(R)=\sum_{m\ge1}w^mR'(m)$, and the polynomial
\[
  T(X)=\nrm\cdot\det\bigl[\mu_X(\varphi(u)\,u^{i+j})\bigr]_{0\le i,j<h}\in\Q[X],
  \qquad \nrm=\frac{(K!/A!^2)^h}{\prod_{i<h}i!^2},
\]
of degree at most $h$. We use two families of parameters $(h,K,A)$, indexed by $n\ge1$:
\begin{center}
\begin{tabular}{@{}lllll@{}}
\toprule
 & $(h,K,A)$ & $\kappa=K/h$ & $\alpha=A/h$ & used for\\
\midrule
general construction $\Gen$ & $(6n,\,7n,\,0)$ & $7/6$ & $0$ & odd $q\ge83$\\
band construction $\Band$ & $(40n,\,49n,\,2n)$ & $49/40$ & $1/20$ & $q=81$\\
\bottomrule
\end{tabular}
\end{center}
For $\Gen$ the only change with respect to Krisman's determinant is that there are $7n$ poles
for $6n$ rows. For $\Band$ the weight also vanishes to order two at $u=1,\dots,A$, which removes
the first $A$ terms of the series. We write $d=\kappa-2\alpha$, so $d=7/6$ for $\Gen$ and
$d=9/8$ for $\Band$.

\subsection{The argument in one page}\label{subsec:one-page}

For a prime $p$ and a nonzero $F\in\Q[X]$ let $\vG p(F)$ be the minimum of the $p$-adic
valuations of the coefficients of $F$. Let $\ell=\prod_pp^{-\vG p(T)}$, so that $\ell T\in\Z[X]$.
If $\xi=a/b$ were rational, then $b^h\ell T(a/b)$ would be an integer. The proof consists in
showing that this integer is nonzero for infinitely many $n$, while its absolute value is at
most $\exp\bigl(h^2(\bconst-\aconst+o(1))\bigr)$ with $\bconst<\aconst$
(\cref{prop:criterion}). This needs three estimates.

\begin{theorem}\label{thm:three-estimates}
Let $\mathcal C\in\{\Gen,\Band\}$, let $q\ge5$ be odd and $\sigma=\pm1$, and let
$T=T_n$ be the polynomial attached to the $n$-th member of $\mathcal C$, with $h=h_n$.
\begin{enumerate}[(i)]
\item \emph{Arithmetic.} Along the $n$ with $T_n\ne0$,
\[
  \liminf\frac{1}{h^2}\sum_p\vG p(T)\log p\ \ge\ \aconst_\sigma(q):=
  \Iprof_{\mathcal C}+c_{2,\mathcal C}\log2+\frac13\log q-\frac{(2-d)^2}4\log(q^2-4\sigma),
\]
with the constants $c_{2,\mathcal C}$ and $\Iprof_{\mathcal C}$ of \cref{tab:constants}.
\item \emph{Size.} For every comparison density $\rho$ and every $\lambda$ satisfying the tail
condition of \cref{thm:decay}, $\limsup h^{-2}\log\abs{T(\xi)}\le\bconst(\rho,\lambda)$, an
explicit number that depends on $q$ only through $W$.
\item \emph{Nonvanishing.} Let $r\in\Q$. If $\mathcal C=\Gen$ and $q\ge83$, or
$\mathcal C=\Band$ and $q=81$, then $T_n(r)\ne0$ for every sufficiently large prime $n$.
\end{enumerate}
\end{theorem}

\begin{proof}
Part~(i) is \cref{thm:arithmetic}, part~(ii) is \cref{thm:decay}, and part~(iii) is
\cref{thm:nonvanishing} together with \cref{prop:local-dets-general,prop:local-dets-band}.
\end{proof}

\begin{table}[ht]
\centering
\caption{The contributions to $\aconst_\sigma(q)$ and the comparison with $\bconst$, per
$h^2$. The constant $\aconst_-$ is the smaller of the two arithmetic constants
(\cref{rem:worse-sign}); $\bconst$ is the certified size exponent
(Appendix~\ref{app:decay-certificates}). Decimals are rounded in the safe direction; the first three rows are evaluated for $\sigma=-1$ and rounded to six places.}
\label{tab:constants}
\begin{tabular}{@{}lll@{}}
\toprule
 & $\Gen$, $q=83$ & $\Band$, $q=81$\\
\midrule
prime $2$ (\cref{thm:prime-two}) & \shortstack[l]{$-\frac{140}{99}\log2$\\$\approx-0.980208$} &
  \shortstack[l]{$-\frac{1847}{1100}\log2$\\$\approx-1.163857$}\\[4pt]
primes $p\mid q$ (\cref{prop:p-divides-q}) & \shortstack[l]{$+\frac13\log q$\\$\approx+1.472947$} &
  \shortstack[l]{$+\frac13\log q$\\$\approx+1.464816$}\\[4pt]
primes $p\mid q^2-4\sigma$ (\cref{prop:p-divides-Y}) &
  \shortstack[l]{$-\frac{25}{144}\log(q^2-4\sigma)$\\$\approx-1.534420$} &
  \shortstack[l]{$-\frac{49}{256}\log(q^2-4\sigma)$\\$\approx-1.682367$}\\[4pt]
ordinary primes, $\Iprof_{\mathcal C}$ (\cref{thm:ordinary}) & $\ge-2.113164137$ &
  $\ge-1.987089582$\\
\midrule
$\aconst_-(q)$ & $\ge-3.154846$ & $\ge-3.368498$\\
$\bconst$ & $\le-3.161708$ & $\le-3.374176$\\
margin $\aconst_--\bconst$ & $\ge0.006862$ & $\ge0.005678$\\
\bottomrule
\end{tabular}
\end{table}

The estimates are obtained as follows.
\begin{itemize}
\item \emph{Prime $2$} (\cref{sec:prime-two}). In the coordinate $z=2u$ the poles are odd and
the summation points even, and a Vandermonde-type minor whose nodes share a parity is divisible by
about $2^{t^2}$; the Cauchy--Binet formula then reduces $\vG2(T)$ to a quadratic minimisation in
the proportion $x$ of pole terms, with minimum at $x=29/33$ for $\Gen$ and $x=51/55$ for $\Band$.
\item \emph{Primes dividing $q(q^2-4\sigma)$} (\cref{sec:exceptional}). For $p\mid q$, reflecting
$m\mapsto-m$ makes the series $p$-adically convergent, and each level offers at most three terms
of valuation at least $2v_p(q)m$, which gives $\frac13v_p(q)h^2$; for $p\mid q^2-4\sigma$ the only
loss comes from the denominators $(1-w)^{b+1}$, and a convexity argument gives
$\frac{(2-d)^2}4v_p(q^2-4\sigma)h^2$.
\item \emph{Ordinary primes} (\cref{sec:ordinary}). For $p\asymp h$ two lower bounds, by Taylor
jets in each residue class modulo $p$ and by global partial fractions with Vandermonde gains, are
small minimisation problems over the residue classes; their maximum, integrated against $\log p$
by the prime number theorem, gives $\Iprof_{\mathcal C}$; the integral is evaluated exactly (Appendix~\ref{app:profile-certificate}).
\item \emph{Size} (\cref{sec:decay}). The determinant at $X=\xi$ is a sum over configurations of
$h$ points on the summation sites; Vandermonde bounds and a discrete energy inequality compare it
with an explicit step density, and the resulting number $\bconst$ is evaluated in ball arithmetic
(Appendix~\ref{app:decay-certificates}).
\item \emph{Nonvanishing} (\cref{sec:nonvanishing}). For $h=tp$ with $p$ prime, a basis adapted
to the residue classes modulo $p$ makes the scaled determinant block triangular modulo $p$, with
diagonal blocks equal, up to units, to three fixed local Hankel determinants, which are nonzero.
\item \emph{Assembly} (\cref{sec:proofs}). The margins of \cref{tab:constants} are positive; a
shift of $\lambda$ (\cref{prop:shift}) carries the case $q=83$ to every odd $q\ge83$. For each $q$
the inequality to be proved is $\bconst_q<\aconst_-(q)$, where $\bconst_q\le\bconst_{83}-\frac12\log\frac q{83}$
and $\aconst_-(q)\ge\aconst_-(83)-\frac1{72}\log\frac q{83}$; this is the whole uniformity
argument in $q$.
\end{itemize}

\subsection{What is new and what is inherited}\label{subsec:new}

The following ingredients come from Krisman's work~\cite{Krisman2026chi2}, in some cases in a
modified form: the functional and its pole series (\cref{lem:pole-series}); the irrationality
criterion (\cref{prop:criterion}); the reflection at primes dividing $q$ and the count of three
terms per level (\cref{prop:p-divides-q}); the degree count at primes dividing $q^2-4\sigma$
(\cref{prop:p-divides-Y}); the size bound by a discrete logarithmic gas compared with an explicit
step density, and its shift to smaller $\abs w$ (\cref{sec:decay}); and the strategy of proving
nonvanishing by reduction modulo a prime proportional to $h$ (\cref{sec:nonvanishing}). We reprove
all of them, because the constructions differ. What is new is the following.
\begin{enumerate}[(a)]
\item \emph{The ratio $K/h>1$.} More poles than rows lowers the degree of the polynomial part of
the weights. At primes dividing $q^2-4\sigma$ the loss drops from $\frac14\log(q^2-4\sigma)$ to
$\frac{25}{144}\log(q^2-4\sigma)$ per $h^2$, which at $q=83$ saves about $0.68$ and outweighs the
larger costs at $2$ and at the ordinary primes. With $K>h$ there is no Lagrange basis at the
poles, so the residue-class and nonvanishing arguments are set up with Taylor jets instead.
\item \emph{The prime $2$.} The parity argument gives the cost $\frac{140}{99}\log2$ for $\Gen$.
For Krisman's ratio $\kappa=1$ the limit of the bound is $\frac{13}{11}\log2\approx1.18\log2$,
which lies between his lower bound $\log2$ and his upper bound $\frac43\log2$
\cite[Lemma~11.1, Prop.~11.2]{Krisman2026chi2}. He also reports a measured loss of about
$1.18\log2$ per $h^2$ at $h=64$; we do not claim agreement with it (\cref{rem:two-comparison} and
Table~\ref{tab:two-adic-square} compare the finite values in two normalisations).
\item \emph{The ordinary primes.} The profile is valid for every $\kappa$ and with zeros in the
numerator. At $\kappa=1$ its integral equals Krisman's generic constant (\cref{rem:profile-k1}),
so it reorganises his class bound rather than improving it.
\item \emph{The zero band} (only for $q=81$). It costs in the size exponent and at the prime $2$,
but changes the balance between the remaining terms enough to reach $q=81$.
\item \emph{Nonvanishing.} The local determinants are $6\times6$ for $\Gen$, explicit
polynomials in $Q=\sigma q^2/4$ without roots of modulus $189$ or more, and $40\times40$ for
$\Band$, computed exactly at $q=81$.
\end{enumerate}

\subsection{Organisation and notation}\label{subsec:notation}

\Cref{sec:pencils} defines the functional, the determinants and the criterion.
\Cref{sec:prime-two,sec:exceptional,sec:ordinary} bound the denominators, \cref{sec:decay} the
size, and \cref{sec:nonvanishing} proves nonvanishing; \cref{sec:proofs} assembles the proof and
records the status of each step, and \cref{sec:limitations} explains why the method stops at $81$.
Appendix~\ref{app:certificates} describes the finite computations and
Appendix~\ref{app:consistency} compares the estimates with exact computations at small sizes.
\Cref{tab:notation} lists the notation used throughout.

\begin{table}[ht]
\centering
\caption{Notation.}\label{tab:notation}
\small
\begin{tabular}{@{}p{0.27\textwidth}p{0.66\textwidth}@{}}
\toprule
$q$, $\sigma$; $y$, $w$, $W$, $\xi$ & odd integer $\ge5$, sign $\pm1$; see~\eqref{eq:basic-notation}\\
$\tau_k$ & partial sum $\sum_{j\le k}\sigma^jy^{-2j-1}/(2j+1)^2$ of $\xi$, \eqref{eq:tau}\\
$h$, $K$, $A$; $\kappa$, $\alpha$, $d$ & size, number of poles, length of the zero band;
$K/h$, $A/h$, $\kappa-2\alpha$\\
$x_k$, $D_K$, $N_A$ & $-k-\frac12$; $\prod_{k<K}(u-x_k)$; $\prod_{j=1}^A(u-j)^2$\\
$\varphi$, $\Phi$ & $uN_A/D_K$; $(K!/A!^2)\,\varphi$, see~\eqref{eq:Phi}\\
$\Omega_b$, $\mom$, $\mu_X$ & $w^{\max(b,1)}(1-w)^{-b-1}$; rational moment sum; the functional
\eqref{eq:functional}\\
$\Delta$, $\nrm$, $T$ & Hankel determinant, normaliser, $T=\nrm\Delta$, see~\eqref{eq:Delta-T}\\
$\beta_i$, $\vvec(u)$, $\bvec(u)$ & $\binom ui$; the columns $(u^i)_{i<h}$ and $(\beta_i(u))_{i<h}$\\
$\vG p$ & Gauss valuation of a polynomial in $X$ over $\Q_p$\\
$\aconst_\sigma(q)$, $\bconst$, $\Iprof_{\mathcal C}$ & arithmetic constant, size exponent,
ordinary-prime constant (\cref{tab:constants})\\
$\psi(t)$, $\pos x$, $s_2(n)$ & $t\log\abs t-t$ with $\psi(0)=0$; $\max(x,0)$; binary digit
sum, so $v_2(n!)=n-s_2(n)$\\
\bottomrule
\end{tabular}
\end{table}

\emph{Conventions.} $p$ is always a prime and $v_p$ the $p$-adic valuation on $\Q_p$. The Gauss
valuation satisfies $\vG p(FG)=\vG p(F)+\vG p(G)$ and $\vG p(F+G)\ge\min(\vG p(F),\vG p(G))$, with
$\vG p(0)=+\infty$, and is invariant under $X\mapsto X+c$ for $c\in\Z_p$; hence the determinant of
a matrix over $\Q_p[X]$ has Gauss valuation at least the minimum over the terms of its Leibniz
expansion. For $u$ in any $\Q$-algebra, $\binom uk=u(u-1)\cdots(u-k+1)/k!$. The implied constants
in $O(\cdot)$ may depend on $q$ and on the construction, but not on $n$.

\section{The half-integer pencils}\label{sec:pencils}

This section defines the objects of the proof: the functional $\mu_X$, the polynomial $T$, its
expansion over the summation points, and the irrationality criterion. All later sections are
statements about $T$, and the expansion~\eqref{eq:real-sites} with its $p$-adic analogues is the
starting point of \cref{sec:prime-two,sec:exceptional,sec:ordinary,sec:decay}.

Throughout, $q\ge5$ is odd, $\sigma=\pm1$, and $y,w,W,\xi$ are as
in~\eqref{eq:basic-notation}. We put
\begin{equation}\label{eq:tau}
  \tau_k=\sum_{j=0}^{k}\frac{\sigma^jy^{-2j-1}}{(2j+1)^2}\qquad(k\ge0),
\end{equation}
the partial sums of the series defining $\xi$.

\subsection{Moment sums and the functional}\label{subsec:moment-sums}

For $b\ge0$ let $\Omega_b(w)=w^{\max(b,1)}(1-w)^{-b-1}\in\Q(w)$. For a polynomial $g\in\Q[u]$,
expanded in the Newton basis as $g(u)=\sum_b\Delta^bg(0)\binom ub$, put
\[
  \mom(g)=\sum_{b\ge0}\Delta^bg(0)\,\Omega_b(w).
\]

\begin{lemma}\label{lem:moment-sums}
Let $g\in\Q[u]$ and $w\in\Q\setminus\{0,1\}$.
\begin{enumerate}[(a)]
\item The identity $\sum_{m\ge1}g(m)w^m=\mom(g)$ holds in $\Q[[w]]$. Consequently, since the expansion of
$w^b(1-w)^{-b-1}$ converges for $\abs w_v<1$ in every completion, for every
absolute value $\abs\cdot_v$ on $\Q$ with $\abs w_v<1$, the series $\sum_{m\ge1}g(m)w^m$
converges in the completion $\Q_v$ to $\mom(g)$.
\item The identity $\mom(g)=-\sum_{m\ge0}g(-m)w^{-m}$ holds in $\Q[[w^{-1}]]$, and hence in
$\Q_v$ whenever $\abs w_v>1$.
\item Let $p$ be a prime such that $w$ and $1-w$ are $p$-adic units. If $g(\Z)\subset\Z_p$, then
$\mom(g)\in\Z_p$.
\end{enumerate}
\end{lemma}

\begin{proof}
(a) For $b\ge1$, $\sum_{m\ge1}\binom mbw^m=w^b(1-w)^{-b-1}$, and $\sum_{m\ge1}w^m=w(1-w)^{-1}$.
(b) We have $\binom{-m}{b}=(-1)^b\binom{m+b-1}{b}$, so for $b\ge1$, with $v=w^{-1}$,
\[
 \sum_{m\ge0}\binom{-m}bv^m=(-1)^b\frac{v}{(1-v)^{b+1}}=(-1)^b\frac{w^{b}}{(w-1)^{b+1}}
 =-\Omega_b(w),
\]
and $\sum_{m\ge0}v^m=w/(w-1)=-\Omega_0(w)$. (c) The Newton coefficients $\Delta^bg(0)$ are
integral combinations of $g(0),\dots,g(b)$, and $\Omega_b(w)\in\Z_p$.
\end{proof}

Let $K\ge1$, $x_k=-k-\frac12$ for $0\le k<K$, and $D_K(u)=\prod_{k<K}(u-x_k)$. Let $\Rsp$ be
the $\Q$-vector space spanned by $\Q[u]$ and the fractions $1/(u-x_k)$, $0\le k<K$; every
rational function whose only poles are simple poles at some of the $x_k$ lies in $\Rsp$, and
its decomposition into a polynomial part and principal parts is unique. For an indeterminate
$X$ define the $\Q$-linear map $\mu_X\colon\Rsp\to\Q[X]$ by
\begin{equation}\label{eq:functional}
  \mu_X(R)=\mom(R')\quad(R\in\Q[u]),\qquad
  \mu_X\Bigl(\frac1{u-x_k}\Bigr)=-4\sigma^ky^{2k+1}\,(X-\tau_k).
\end{equation}
Each value is a polynomial of degree at most one in $X$.

\begin{lemma}[pole series]\label{lem:pole-series}
For every $R\in\Rsp$ the series $\sum_{m\ge1}w^mR'(m)$ converges absolutely and equals
$\mu_\xi(R)$.
\end{lemma}

\begin{proof}
For polynomials this is \cref{lem:moment-sums}(a) with the real absolute value. For
$R=1/(u-x_k)$ we have $R'(m)=-4/(2m+2k+1)^2$. Writing $2m+2k+1=2j+1$ with $j=m+k\ge k+1$
and $w^m=\sigma^{j-k}y^{-2(j-k)}$,
\[
  \sum_{m\ge1}w^mR'(m)=-4\sigma^ky^{2k+1}\sum_{j>k}\frac{\sigma^jy^{-2j-1}}{(2j+1)^2}
  =-4\sigma^ky^{2k+1}(\xi-\tau_k).\qedhere
\]
\end{proof}

The same computation will be used $2$-adically in \cref{sec:prime-two} and $p$-adically in
\cref{sec:exceptional}. Lemma~\ref{lem:pole-series} is Krisman's pole series
\cite[Lemma~2.1]{Krisman2026chi2}, with an arbitrary number $K$ of poles.

\subsection{The determinants}\label{subsec:determinants}

Let $h\ge1$, $K\ge h$ and $A\ge0$ be integers, and put
\begin{equation}\label{eq:phi}
  N_A(u)=\prod_{j=1}^A(u-j)^2,\qquad \varphi(u)=\frac{u\,N_A(u)}{D_K(u)}.
\end{equation}
For every polynomial $g$ the product $\varphi g$ lies in $\Rsp$. Define
\begin{equation}\label{eq:Delta-T}
  \Delta(X)=\det\bigl[\mu_X(\varphi(u)u^{i+j})\bigr]_{0\le i,j<h},\qquad
  \nrm=\frac{(K!/A!^2)^h}{\prod_{i<h}i!^2},\qquad T(X)=\nrm\,\Delta(X).
\end{equation}
These are polynomials in $X$ with rational coefficients and degree at most $h$, and
$\Delta(\xi)$ is the determinant of the real matrix
$\bigl[\sum_{m\ge1}w^m(\varphi u^{i+j})'(m)\bigr]$ by \cref{lem:pole-series}.

\begin{lemma}[change of basis]\label{lem:change-of-basis}
Let $a_0,\dots,a_{h-1}$ and $b_0,\dots,b_{h-1}$ be polynomials of degree less than $h$ with
coefficient matrices $C_a$, $C_b$ in the monomial basis. Then
$\det[\mu_X(\varphi a_ib_j)]=\det C_a\det C_b\,\Delta(X)$. In particular, with
$\beta_i(u)=\binom ui$ and
\begin{equation}\label{eq:Phi}
  \Phi(u)=\frac{K!}{A!^2}\varphi(u)=\frac{K!\,u\binom{u-1}{A}^2}{D_K(u)},
\end{equation}
one has $T(X)=\det[\mu_X(\Phi\beta_i\beta_j)]_{0\le i,j<h}$.
\end{lemma}

\begin{proof}
The matrix $[\mu_X(\varphi a_ib_j)]$ equals $C_a[\mu_X(\varphi u^{r+s})]_{r,s}C_b^{\mathsf T}$
by bilinearity. For the second statement, the matrix of $(\beta_i)$ is lower triangular with
diagonal $1/i!$, and $N_A(u)=A!^2\binom{u-1}{A}^2$.
\end{proof}

\emph{The expansion over summation points.} Let $\vvec(u)=(u^i)_{0\le i<h}$ be the column of monomials and $\vvec'(u)$ its derivative. For
$X=\xi$, \cref{lem:pole-series} and $(u^{i+j})'=(u^i)'u^j+u^i(u^j)'$ give
\begin{equation}\label{eq:real-sites}
  \bigl[\mu_\xi(\varphi u^{i+j})\bigr]_{i,j}
  =\sum_{m>A}w^m\,\bigl[\vvec(m)\ \vvec'(m)\bigr]
  \begin{pmatrix}\varphi'(m)&\varphi(m)\\ \varphi(m)&0\end{pmatrix}
  \bigl[\vvec(m)\ \vvec'(m)\bigr]^{\mathsf T},
\end{equation}
an absolutely convergent sum of matrices of rank at most two; the sites $1\le m\le A$ do not
contribute because $\varphi$ and $\varphi'$ vanish there. We call the two columns
$\vvec(m)$ and $\vvec'(m)$ the \emph{slots} of the site $m$. The same structure, with other
sites and with additional rank-one terms at the poles, appears $2$-adically and $p$-adically in
\cref{sec:prime-two,sec:exceptional}. All of these expansions are used through the
Cauchy--Binet formula: if $M=EBE^{\mathsf T}$ with $E$ of size $h\times N$ and $B$ of size
$N\times N$, then
\begin{equation}\label{eq:cauchy-binet}
  \det M=\sum_{S,S'}\det E_{[h],S}\,\det B_{S,S'}\,\det E_{[h],S'},
\end{equation}
the sum running over pairs of $h$-element subsets of the column indices of $E$. When $B$ is
block diagonal, $\det B_{S,S'}=0$ unless $S$ and $S'$ meet every block in the same number of
indices, and then it is the product of the corresponding minors of the blocks.

\subsection{The criterion}\label{subsec:criterion}

\begin{proposition}\label{prop:criterion}
Let $\xi\in\R$ and let $(T_n)_{n\in\mathcal N}$, $\mathcal N\subset\N$ infinite, be nonzero
polynomials in $\Q[X]$ with $\deg T_n\le h_n$, $h_n\to\infty$. Suppose that, along $\mathcal N$,
\[
  \liminf\frac1{h_n^2}\sum_p\vG p(T_n)\log p\ \ge\ \aconst,\qquad
  \limsup\frac1{h_n^2}\log\abs{T_n(\xi)}\ \le\ \bconst,
\]
with $\bconst<\aconst$, and that for every $r\in\Q$ there are infinitely many $n\in\mathcal N$
with $T_n(r)\ne0$. Then $\xi\notin\Q$.
\end{proposition}

\begin{proof}
Let $\ell_n=\prod_pp^{-\vG p(T_n)}$; then $\ell_nT_n\in\Z[X]$. If $\xi=a/b$ with $b\ge1$, then
$b^{h_n}\ell_nT_n(a/b)$ is an integer, nonzero for infinitely many $n$, while
$\log\abs{b^{h_n}\ell_nT_n(\xi)}\le h_n\log b+(\bconst-\aconst+o(1))h_n^2\to-\infty$.
\end{proof}

This is Krisman's criterion \cite[Prop.~3.1, Thm.~3.2]{Krisman2026chi2} in the form we need;
its elementary core is formalised in Lean (see the section on reproducibility). We apply it with
$\mathcal N$ the set of all $n$ with $T_n\ne0$ for the arithmetic and analytic estimates, and with
the prime values of $n$ for nonvanishing.

\section{The prime 2}\label{sec:prime-two}

This section proves the lower bound for $\vG2(T)$ that gives the term $c_{2,\mathcal C}\log2$
of $\aconst_\sigma(q)$ (\cref{thm:prime-two}); it is used in \cref{thm:arithmetic}. The prime $2$
needs its own argument because $v_2(y)=-1$: the poles $x_k$ are not $2$-adically integral. In the
coordinate $z=2u$ they become the odd integers $z_k=-2k-1$, while the summation points become the
even integers $2m$, and the gain comes from treating the two parities separately.

\subsection{The rescaled determinant and its \texorpdfstring{$2$}{2}-adic expansion}\label{subsec:rescaled}

Put $P_K(z)=\prod_{k<K}(z-z_k)$, $N^{(2)}(z)=\prod_{j=1}^A(z-2j)^2$ and $F(z)=zN^{(2)}(z)/P_K(z)$.
Since $D_K(u)=2^{-K}P_K(2u)$ and $N_A(u)=2^{-2A}N^{(2)}(2u)$,
\begin{equation}\label{eq:phi-F}
  \varphi(u)=2^{K-1-2A}F(2u).
\end{equation}
On the space spanned by $\Q[z]$ and the fractions $1/(z-z_k)$ define the affine functional
\[
  \nu_X(G)=\mom\bigl(m\mapsto G'(2m)\bigr)\ (G\in\Q[z]),\qquad
  \nu_X\Bigl(\frac1{z-z_k}\Bigr)=-\sigma^ky^{2k+1}(X-\tau_k).
\]
If $R(u)=G(2u)$ then $R'(m)=2G'(2m)$, and $1/(u-x_k)=2/(2u-z_k)$; comparing
with~\eqref{eq:functional} gives $\mu_X(R)=2\nu_X(G)$ on both kinds of generators. With
$H(X)=[\nu_X(Fz^{i+j})]_{0\le i,j<h}$ and~\eqref{eq:phi-F}, the $(i,j)$ entry of the matrix
defining $\Delta$ is $2^{K-2A-i-j}H_{ij}$, hence
\begin{equation}\label{eq:Delta-H}
  \Delta(X)=2^{(K-2A)h-h(h-1)}\det H(X).
\end{equation}

The series $\xi_2=\sum_{j\ge0}\sigma^jy^{-2j-1}/(2j+1)^2$ converges in $\Z_2$, because the $j$-th
term has valuation $2j+1$. Repeating the computation in the proof of \cref{lem:pole-series} in
$\Q_2$, where $\sum_mw^m(\cdots)$ converges because $v_2(w)=2$, gives
\[
  \sum_{m\ge1}w^m\frac{\mathrm d}{\mathrm dz}\Bigl(\frac1{z-z_k}\Bigr)\Big|_{z=2m}
  =-\sigma^ky^{2k+1}(\xi_2-\tau_k).
\]
Hence $\nu_X(1/(z-z_k))=\sum_mw^m(\cdots)-\sigma^ky^{2k+1}(X-\xi_2)$. Writing
$Fz^s=Q_s+\sum_kr_kz_k^s/(z-z_k)$ with $Q_s$ a polynomial and $r_k=\Res_{z=z_k}F$, and using
\cref{lem:moment-sums}(a) $2$-adically for $Q_s$, we obtain the identity in $\Q_2[X]$
\begin{multline}\label{eq:H-two-adic}
  H(X)=\sum_{m\ge1}w^m\bigl[\vvec(2m)\ \vvec'(2m)\bigr]
  \begin{pmatrix}F'(2m)&F(2m)\\F(2m)&0\end{pmatrix}
  \bigl[\vvec(2m)\ \vvec'(2m)\bigr]^{\mathsf T}\\
  +(X-\xi_2)\sum_{k<K}\gamma_k\,\vvec(z_k)\vvec(z_k)^{\mathsf T},
\end{multline}
with $\gamma_k=-\sigma^ky^{2k+1}r_k$ and the series converging entrywise in $\Q_2[X]$ for the
Gauss norm. The even sites $2m$ with $1\le m\le A$ contribute nothing, since $F$ and $F'$
vanish there.

The weights have the following valuations. Since $z_k-z_l=2(l-k)$ and $z_k-2j$ is odd,
\begin{equation}\label{eq:gamma-two}
  v_2(\gamma_k)=-(2k+1)-(K-1)-v_2(k!)-v_2((K-1-k)!)\ \ge\ -2k-2K+1 .
\end{equation}
(The inequality holds because $v_2(k!)+v_2((K-1-k)!)\le v_2((K-1)!)\le K-2$, which gives even
$-2k-2K+2$; we keep the weaker form, which is all that is used below.)
For $m>A$ the integer $P_K(2m)$ is odd and
$N^{(2)}(2m)=2^{2A}\bigl(A!\binom{m-1}A\bigr)^2$, so
\[
  v_2(F(2m))\ge1+2A+2v_2(A!)\ge 4A-2\floor{\log_2(A+1)}-1 ,
\]
because $v_2(A!)=A-s_2(A)$ and $s_2(A)\le\floor{\log_2(A+1)}$ (the true bound is
$4A+1-2\floor{\log_2(A+1)}$; we keep the weaker form).
The logarithmic derivative $F'/F=1/z+\sum_j2/(z-2j)-\sum_k1/(z-z_k)$ has valuation at least
$-1-\floor{\log_2m}$ at $z=2m$. Consequently
\begin{equation}\label{eq:nu-m}
  \nu_m:=\min\bigl(v_2F(2m),\,v_2F'(2m)\bigr)\ \ge\ 4A-2\floor{\log_2(A+1)}-2-\floor{\log_2 m}
  \qquad(m>A),
\end{equation}
and $\nu_m\ge0$ when $A=0$, because then $F(2m)$ and $F'(2m)$ are $2$-integral.

\subsection{Minors with nodes of one parity}\label{subsec:parity-minors}

For $t\ge0$ put
\[
  \Pa(t)=\sum_{j<t}\bigl(j+v_2(j!)\bigr)=t(t-1)-\sum_{j<t}s_2(j)=t^2+O(t\log(t+2)).
\]

\begin{lemma}[parity minors]\label{lem:parity-minors}
Let $h\ge2$, $L=\floor{\log_2(h-1)}$, $\varepsilon\in\{0,1\}$, and let $J\subset\{0,\dots,h-1\}$
have $t$ elements. Consider a $t\times t$ matrix whose rows are indexed by $J$ and whose columns
are of the form $(\zeta^i)_{i\in J}$ or $(i\zeta^{i-1})_{i\in J}$ with
$\zeta\in\varepsilon+2\Z$, and let $r$ be the number of columns of the second kind. Then its
determinant has $2$-adic valuation at least $\Pa(t)-r(1+L)$.
\end{lemma}

\begin{proof}
Let $B_j(z)=\prod_{i<j}(z-\varepsilon-2i)$, a monic polynomial in $\Z[z]$ of degree $j$. Then
$z^i=\sum_{j\le i}c_{ij}B_j(z)$ with $c_{ij}\in\Z$ and $c_{ii}=1$. For $\zeta=2s+\varepsilon$,
$s\in\Z$, we have $B_j(\zeta)=2^jj!\binom sj$, which is divisible by $2^{j+v_2(j!)}$. For the
derivative, the identity $\frac{\mathrm d}{\mathrm ds}\binom sj=\sum_{l=1}^{j}
\frac{(-1)^{l-1}}{l}\binom s{j-l}$ (compare coefficients of $T^j$ in
$\partial_s(1+T)^s=\log(1+T)(1+T)^s$) gives $B_j'(\zeta)=2^{j-1}j!\,\partial_s\binom sj$, of
valuation at least $j+v_2(j!)-1-\floor{\log_2j}$. By the Cauchy--Binet formula for the product
of the $t\times h$ matrix $(c_{ij})_{i\in J,\,j<h}$ and the $h\times t$ matrix of the values
$B_j(\zeta)$ or $B'_j(\zeta)$, the determinant is an integral combination of minors with rows
$j_1<\dots<j_t$ of the latter matrix. Since $j\mapsto j+v_2(j!)$ is increasing and
$j_r\ge r-1$, each such minor has valuation at least $\sum_{i<t}(i+v_2(i!))-r(1+L)$.
\end{proof}

\begin{corollary}\label{cor:parity-laplace}
Let $\mathsf E$ be an $h\times h$ matrix whose columns are $P$ columns $\vvec(\zeta)$ with $\zeta$ odd
and $h-P$ columns $\vvec(\zeta)$ or $\vvec'(\zeta)$ with $\zeta$ even, $r$ of them derivatives.
Then $v_2(\det\mathsf E)\ge\Pa(P)+\Pa(h-P)-r(1+L)$.
\end{corollary}

\begin{proof}
Laplace expansion along the odd columns writes $\det\mathsf E$ as a signed sum, over $P$-element sets
$J$ of rows, of products of a $P\times P$ minor with odd nodes and rows $J$ and an
$(h-P)\times(h-P)$ minor with even nodes and the complementary rows. Apply
\cref{lem:parity-minors} to both.
\end{proof}

For nodes of mixed parity the corresponding bound would be the valuation of a Vandermonde
determinant of $t$ consecutive integers, about $t^2/2$; the separation by parity doubles it,
and this is where the constant of \cref{thm:prime-two} comes from.

\subsection{The bound at the prime 2}\label{subsec:two-bound}

\begin{theorem}\label{thm:prime-two}
Let $q$ be odd, $\sigma=\pm1$, and $K\ge h\ge2$, $A\ge0$. Let $\lambda_0\le\lambda_1\le\cdots\le\lambda_{K-1}$ be the exact valuations $v_2(\gamma_k)$, $k<K$, of~\eqref{eq:gamma-two} in increasing order, and $\epsilon_1\le\epsilon_2\le\cdots$ the numbers $2m+\nu_m$, $m>A$, each listed twice, in increasing order. Then
\begin{multline*}
  \vG2(T)\ \ge\ v_2(\nrm)+(K-2A)h-h(h-1)\\
  +\min_{0\le P\le h}\Bigl[2\Pa(P)+2\Pa(h-P)-(h-P)(1+L)+\sum_{i<P}\lambda_i
  +\sum_{i=1}^{h-P}\epsilon_i\Bigr].
\end{multline*} If $K=\kappa h$ and $A=\alpha h$ with $\kappa\ge1$,
$0\le\alpha$ and $d=\kappa-2\alpha>0$ fixed, then
\begin{gather}
  \vG2(T)\ \ge\ h^2\min_{0\le x\le1}f(x)-O(h\log h),\label{eq:two-asymptotic}\\
  f(x)=2(d-1)-4\kappa x+3x^2+\tfrac52(1-x)^2+6\alpha(1-x).\notag
\end{gather}
For $\Gen$ the minimum is $-140/99$, attained at $x=29/33$; for $\Band$ it is $-1847/1100$,
attained at $x=51/55$.
\end{theorem}

\begin{proof}
By~\eqref{eq:Delta-H}, $\vG2(T)=v_2(\nrm)+(K-2A)h-h(h-1)+\vG2(\det H)$. Truncate the sum over
$m$ in~\eqref{eq:H-two-adic} at $m\le M$; the truncated matrices converge to $H$ in the Gauss
norm, so it suffices to bound the determinant of each truncation uniformly in $M$. Write the
truncation as $\mathsf E\mathsf B\mathsf E^{\mathsf T}$, where the columns of $\mathsf E$ are the slots $\vvec(2m)$,
$\vvec'(2m)$ ($A<m\le M$) and $\vvec(z_k)$ ($k<K$), and $\mathsf B$ is block diagonal with the $2\times2$
blocks $w^m\bigl(\begin{smallmatrix}F'&F\\F&0\end{smallmatrix}\bigr)(2m)$ and the $1\times1$
blocks $(X-\xi_2)\gamma_k$. By~\eqref{eq:cauchy-binet}, $\det H$ is a sum of terms indexed by
pairs $(S,S')$ of sets of $h$ slots that select the same poles and the same number $n_m\in\{1,2\}$
of slots at each selected site $2m$.

Fix such a term with $P$ poles and $e=h-P$ even slots. The middle factor has Gauss valuation at
least $\sum_{k\in S}v_2(\gamma_k)+\sum_m n_m(2m+\nu_m)$, since a $1\times1$ minor of the block at
$2m$ is $w^mF'(2m)$ or $w^mF(2m)$ and the $2\times2$ minor is $-w^{2m}F(2m)^2$; and
$\vG2(X-\xi_2)=0$. This is at least $\sum_{i<P}\lambda_i+\sum_{i\le e}\epsilon_i$. The
derivative--derivative entry of each block vanishes, so a site carrying one slot on each side
contributes at most one derivative column to the two outer factors together, and a site carrying
two slots contributes two; the total number of derivative columns on both sides is therefore at
most $e$. By \cref{cor:parity-laplace} the two outer factors together have valuation at least
$2\Pa(P)+2\Pa(e)-e(1+L)$. This proves the first inequality.

For the asymptotic form, $v_2(n!)=n-s_2(n)$ gives $v_2(\nrm)=(K-2A)h-h^2+O(h\log h)$ and
$\Pa(t)=t^2+O(h\log h)$ for $t\le h$. The $P$ smallest of the lower
bounds~\eqref{eq:gamma-two} are those with $k=K-1,\dots,K-P$, whose sum is $-4KP+P^2+2P$. For
the even slots, \eqref{eq:nu-m} and $\floor{\log_2m}\le\log_2(2h)+m/h$ give
$2m+\nu_m\ge(2-1/h)m+4A-O(\log h)$; since each site carries at most two slots and the sites
start at $A+1$, the $e$ smallest values sum to at least
$(2-1/h)\sum_{i=1}^{e}(A+\lceil i/2\rceil)+4Ae-O(h\log h)=6Ae+e^2/2-O(h\log h)$. Collecting
terms with $P=xh$, $e=(1-x)h$ gives~\eqref{eq:two-asymptotic}. Finally
$f'(x)=11x-4\kappa-5-6\alpha$ vanishes at $x=(4\kappa+5+6\alpha)/11$, which is $29/33$ for
$(\kappa,\alpha)=(7/6,0)$ and $51/55$ for $(49/40,1/20)$; substituting gives the two values.
\end{proof}

\begin{remark}\label{rem:two-comparison}
For $(\kappa,\alpha)=(1,0)$, Krisman's determinant, \eqref{eq:two-asymptotic} gives
$-\frac{13}{11}h^2$ at $x=9/11$. Krisman proves the bound $-\frac43h^2$
\cite[Lemma~11.1]{Krisman2026chi2}, shows that $-h^2+h$ is attained by the leading coefficient
\cite[Prop.~11.2]{Krisman2026chi2}, and reports a measured $2$-adic loss of about $1.18\log2$ per $h^2$ at $h=64$, in a
normalisation we cannot identify from his text. In ours, the exact values for $q=83$,
$\sigma=1$, $K=h$, $A=0$ are given in \cref{tab:two-adic-square}: the loss of the un-normalised
determinant $\Delta=\det[\mu_X(\varphi u^{i+j})]$ is $1.1758$ per $h^2$ at $h=64$, and that of
the normalised polynomial $T$ of this paper is between $1.00$ and $1.11$ for $24\le h\le64$,
against $\frac{13}{11}\approx1.18$ for the limit of the bound. These are different quantities
and the finite bound is still far from its limit at these sizes, so we draw no conclusion from the
proximity of $1.1758$ to $1.18$. The terms with $P=h$ are the only ones that contribute to the
leading coefficient in $X$, and $f(1)=1-2\kappa-4\alpha$ equals $-1$ at $(\kappa,\alpha)=(1,0)$,
in accordance with his Proposition~11.2. The minimum
of $f$ is attained at an interior $x$, that is, by coefficients of lower degree in $X$. For
$\Gen$ at $q=83$ and $h=6,12,\dots,72$, the finite bound of \cref{thm:prime-two} lies below the
exact value of $\vG2(T)$ by between $3$ and $56$ (\cref{tab:two-adic}); the ratios
$\vG2(T)/h^2$ are not monotone in $h$ and all exceed $-140/99$. This is consistency evidence
only.
\end{remark}

\section{Odd primes dividing \texorpdfstring{$q(q^2-4\sigma)$}{q(q²-4σ)}}\label{sec:exceptional}

This section bounds $\vG p(T)$ for the finitely many odd primes dividing $q(q^2-4\sigma)$, which
give the terms $\frac13\log q$ and $-\frac{(2-d)^2}{4}\log(q^2-4\sigma)$ of $\aconst_\sigma(q)$,
and proves a crude uniform bound for the other odd primes, used for small primes in
\cref{sec:ordinary}. All three rest on one lemma on the partial fractions of the normalised
weights (\cref{lem:newton-quotient}).

Throughout, $p$ is an odd prime. Since $1-w=(q^2-4\sigma)/q^2$, the odd primes fall into three
classes: those dividing $q$ (where $\abs w_p>1$), those dividing $q^2-4\sigma$ (where $w$ is a unit
but $1-w$ is not), and the others, called \emph{ordinary}, for which $w$ and $1-w$ are units. As
$q$ is odd, $q^2-4\sigma$ is odd and prime to $q$. We work with the binomial form
$T(X)=\det[\mu_X(\Phi\beta_i\beta_j)]$ of \cref{lem:change-of-basis}, with $\Phi$ as
in~\eqref{eq:Phi}.

\subsection{Newton quotients}\label{subsec:newton}

\begin{lemma}\label{lem:newton-quotient}
Let $p$ be odd and let $B\in\Q[u]$ have degree at most $D$ and satisfy $B(\Z)\subset\Z_p$. Then
\[
  \frac{K!\,B(u)}{D_K(u)}=Q(u)+\sum_{k<K}\frac{\rho_k}{u-x_k}
\]
with a polynomial $Q$ and
\begin{enumerate}[(a)]
\item $\rho_k=(-1)^kK\binom{K-1}{k}B(x_k)\in\Z_p$;
\item $Q'=0$ if $D\le K$; otherwise $\deg Q'\le D-K-1$ and $p^{\lambda}Q'(\Z)\subset\Z_p$
with $\lambda=\floor{\log_pD}+\floor{\log_p(D-K)}$.
\end{enumerate}
\end{lemma}

\begin{proof}
A polynomial with $p$-integral values on $\Z$ is a $\Z_p$-linear combination of the $\binom un$
and therefore maps $\Z_p$ into $\Z_p$; as $p$ is odd, $\frac12\in\Z_p$. Hence the Newton
coefficients $b_t$ in $B(u)=\sum_tb_t\binom{u+K-\frac12}{t}$, which are finite differences of
$B$ at $\frac12-K$, lie in $\Z_p$, and $B(x_k)\in\Z_p$. Part~(a) follows from
$D_K'(x_k)=\prod_{l\ne k}(l-k)=(-1)^kk!\,(K-1-k)!$. For $t\ge K$ the first $K$ factors of
$t!\binom{u+K-1/2}{t}$ are those of $D_K(u)$, so
\[
  \frac{K!}{D_K(u)}\binom{u+K-\frac12}{t}=\binom{t}{K}^{-1}\binom{u-\frac12}{t-K},
\]
while for $t<K$ the quotient is a proper rational function. Thus
\[
  Q(u)=\sum_{t=K}^{D}b_t\binom tK^{-1}\binom{u-1/2}{t-K}.
\]
By Kummer's theorem
$v_p\binom tK\le\floor{\log_pt}$, and
$\frac{\mathrm d}{\mathrm du}\binom{u-1/2}{s}=\sum_{l=1}^{s}\frac{(-1)^{l-1}}{l}\binom{u-1/2}{s-l}$
(the identity used in the proof of \cref{lem:parity-minors}); the values of $\binom{u-1/2}{r}$
at integers lie in $\Z_p$. This gives~(b).
\end{proof}

We apply the lemma to $B_{ij}(u)=u\binom{u-1}{A}^2\beta_i(u)\beta_j(u)$, of degree
$D_{ij}=i+j+2A+1$, for which $K!B_{ij}/D_K=\Phi\beta_i\beta_j$. Writing $Q_{ij}$ and
$\rho_{ij,k}$ for the corresponding quotient and residues,
\begin{equation}\label{eq:entry-decomposition}
  \mu_X(\Phi\beta_i\beta_j)=\mom(Q_{ij}')-4\sum_{k<K}\rho_{ij,k}\,\sigma^ky^{2k+1}(X-\tau_k).
\end{equation}
For both constructions $D_{ij}\le2h+2A-1<3h$ and $2K-1<3h$; we put $L_p=\floor{\log_p(3h)}$.

\begin{proposition}[ordinary odd primes, uniform bound]\label{prop:uniform}
If $p\nmid 2q(q^2-4\sigma)$, then $\vG p(T)\ge-2hL_p$.
\end{proposition}

\begin{proof}
Here $w$ and $1-w$ are $p$-adic units. We bound the two parts of~\eqref{eq:entry-decomposition}.
By \cref{lem:newton-quotient}(b) and
\cref{lem:moment-sums}(c) applied to $p^{2L_p}Q_{ij}'$, $\mom(Q_{ij}')\in p^{-2L_p}\Z_p$. Since
$p\nmid y$, the polynomial $y^{2k+1}(X-\tau_k)$ has coefficients $y^{2k+1}$ and
$\sum_{j\le k}\sigma^jy^{2k-2j}/(2j+1)^2$, of valuation at least $-2\floor{\log_p(2K-1)}\ge-2L_p$.
With \cref{lem:newton-quotient}(a), every entry has Gauss valuation at least $-2L_p$, and the
Leibniz expansion gives the claim.
\end{proof}

This is the analogue, in the binomial normalisation, of \cite[Lemma~4.2]{Krisman2026chi2}.

\subsection{Primes dividing \texorpdfstring{$q^2-4\sigma$}{q²-4σ}}\label{subsec:p-divides-Y}

\begin{proposition}\label{prop:p-divides-Y}
Let $p\mid q^2-4\sigma$ and $f=v_p(q^2-4\sigma)$. Then
\[
  \vG p(T)\ \ge\ -f\sum_{i<h}\pos{2i+2A+1-K}-2hL_p .
\]
If $K=\kappa h$, $A=\alpha h$ with $0<d=\kappa-2\alpha\le2$, the right-hand side is
$-f\frac{(2-d)^2}{4}h^2-O(h\log h)$; the factor is $25/144$ for $\Gen$ and $49/256$ for $\Band$.
\end{proposition}

\begin{proof}
Now $w$ is a unit and $v_p(1-w)=f$, so $v_p(\Omega_b(w))=-f(b+1)$. The Newton expansion of
$Q_{ij}'$ involves only $b\le D_{ij}-K-1$, and its coefficients lie in $p^{-2L_p}\Z_p$, so
$v_p(\mom(Q_{ij}'))\ge-2L_p-f\pos{i+j+2A+1-K}$; the pole terms are as in
\cref{prop:uniform}. In the Leibniz expansion, for every permutation $\pi$, convexity of
$t\mapsto\pos t$ gives
\begin{align*}
  \sum_i\pos{i+\pi(i)+2A+1-K}
  &\le\sum_i\tfrac12\bigl(\pos{2i+2A+1-K}+\pos{2\pi(i)+2A+1-K}\bigr)\\
  &=\sum_i\pos{2i+2A+1-K}.
\end{align*}
Finally $\sum_{i<h}\pos{2i+1-dh}=h^2\int_0^1\pos{2t-d}\diff t+O(h)=\frac{(2-d)^2}4h^2+O(h)$.
\end{proof}

At $\kappa=1$, $\alpha=0$ this is the bound $\frac14h^2\log(q^2-4)$ of
\cite[Lemma~5.1, Thm.~5.3]{Krisman2026chi2}. The decrease from $\frac14$ to $\frac{25}{144}$
is the main gain of the ratio $\kappa=7/6$.

\subsection{Primes dividing \texorpdfstring{$q$}{q}: reflection}\label{subsec:p-divides-q}

Let $p\mid q$ and $f=v_p(q)\ge1$; then $v_p(y)=f$ and $v_p(w)=-2f$. The series
\[
  \xi'_p=\sum_{j\ge0}\frac{\sigma^jy^{2j+1}}{(2j+1)^2}
\]
converges in $\Z_p$: its $j$-th term has valuation $(2j+1)f-2v_p(2j+1)\ge(2j+1)-2\log_3(2j+1)>0$,
which tends to infinity.

\begin{lemma}[reflected functional]\label{lem:reflection}
For every $R\in\Rsp$,
\[
  \mu_X(R)=-\sum_{m\ge0}w^{-m}R'(-m)\;-\;4\,(X+\sigma\xi'_p)\sum_{k<K}\sigma^ky^{2k+1}
  \Res_{u=x_k}R,
\]
where the series converges in $\Q_p$.
\end{lemma}

\begin{proof}
For polynomials this is \cref{lem:moment-sums}(b), since $\abs{w^{-1}}_p<1$. For
$R=1/(u-x_k)$, $R'(-m)=-4/(2k+1-2m)^2$. The terms $0\le m\le k$ give, with $j=k-m$,
$4\sum_{j\le k}\sigma^{k-j}y^{2k-2j}/(2j+1)^2=4\sigma^ky^{2k+1}\tau_k$, and the terms $m>k$
give, with $j=m-k-1$, $4\sigma^{k+1}y^{2k+1}\xi'_p$. Hence
$-\sum_mw^{-m}R'(-m)=4\sigma^ky^{2k+1}(\tau_k+\sigma\xi'_p)$, and adding
$-4\sigma^ky^{2k+1}(X+\sigma\xi'_p)$ gives $-4\sigma^ky^{2k+1}(X-\tau_k)=\mu_X(R)$.
\end{proof}

\begin{proposition}\label{prop:p-divides-q}
Let $p\mid q$ and $f=v_p(q)$. Then $\vG p(T)\ge\frac f3h^2-O(h\log h)$.
\end{proposition}

\begin{proof}
Let $\bvec(u)=(\beta_i(u))_{i<h}$. By \cref{lem:reflection} applied to $\Phi\beta_i\beta_j$,
\begin{multline*}
  \bigl[\mu_X(\Phi\beta_i\beta_j)\bigr]=\sum_{m\ge0}w^{-m}\bigl[\bvec\ \bvec'\bigr](-m)
  \begin{pmatrix}-\Phi'(-m)&-\Phi(-m)\\-\Phi(-m)&0\end{pmatrix}
  \bigl[\bvec\ \bvec'\bigr](-m)^{\mathsf T}\\
  +\sum_{k<K}\eta_k\,\bvec(x_k)\bvec(x_k)^{\mathsf T},
\end{multline*}
with $\eta_k=-4(X+\sigma\xi'_p)\sigma^ky^{2k+1}\rho_k$ and
$\rho_k=\Res_{x_k}\Phi=(-1)^kK\binom{K-1}kx_k\binom{x_k-1}{A}^2\in\Z_p$
(\cref{lem:newton-quotient}(a)). We record the valuations.
\begin{itemize}
\item $\bvec(-m)$ and $\bvec(x_k)$ are integral, and $\beta_i'=\sum_{l=1}^{i}\frac{(-1)^{l-1}}{l}
\beta_{i-l}$ shows that $\bvec'(-m)$ has valuation at least $-\floor{\log_ph}$.
\item $\vG p(\eta_k)\ge(2k+1)f$, because $X+\sigma\xi'_p$ has Gauss valuation $0$.
\item $v_p(w^{-m})=2fm$.
\item $2^KD_K(-m)=\prod_{k<K}(2k+1-2m)$ is a product of $K$ consecutive odd integers, none zero.
For each $e\ge1$ at most $\floor{K/p^e}+1$ of them are divisible by $p^e$, and none if
$p^e>2(K+m)$. Hence $v_p(K!/D_K(-m))\ge-\floor{\log_p(2K+2m)}$. Differentiating
$\Phi=K!\,u\binom{u-1}A^2/D_K$ by the product rule gives
$\Phi'=\frac{K!}{D_K}\bigl[\binom{u-1}A^2+2u\binom{u-1}A\frac{\mathrm d}{\mathrm du}\binom{u-1}A
-u\binom{u-1}A^2\frac{D_K'}{D_K}\bigr]$. At $u=-m$ the binomials are integers, the derivative
$\frac{\mathrm d}{\mathrm du}\binom{u-1}A$ has valuation at least $-\floor{\log_pA}$ by the
identity used in \cref{lem:parity-minors}, and
$D_K'/D_K(-m)=\sum_k2/(2k+1-2m)$ has valuation at least $-\floor{\log_p(2K+2m)}$. Hence
$v_p(\Phi(-m))$, $v_p(\Phi'(-m))\ge-C\log(h+m+2)$ for a constant $C=C(p)$ that depends on $p$
and on the construction.
\end{itemize}
Truncate the sum at $m\le M$ and apply~\eqref{eq:cauchy-binet} with the slots $\bvec(-m)$,
$\bvec'(-m)$ and $\bvec(x_k)$. A nonzero term selects at each site $-m$ the same number of slots
on both sides and the same poles. Its outer factors have valuation at least $-2h\floor{\log_ph}$,
and its middle factor at least the sum, over the $h$ selected slots, of $2fm-C\log(h+m+2)$ for a
slot at $-m$ and of $(2k+1)f\ge2fk$ for the pole $x_k$. Every level $m\ge0$ thus offers at most
three slots, each of cost at least
\[
  2fm-C\log(h+m+2)\ \ge\ \bigl(2f-C/(h+2)\bigr)m-C\log(h+2),
\]
and, for $h$ so large that $2f>C/(h+2)$ (which keeps the costs increasing in $m$), the sum of
the $h$ smallest such costs is at least
$(2f-C/(h+2))\sum_{i<h}\floor{i/3}-Ch\log(h+2)=\frac f3h^2-O(h\log h)$, uniformly in $M$. Since
$w^{-m}\to0$ in $\Q_p$, the truncated determinants converge to $T$ in the Gauss norm.
\end{proof}

This is Krisman's reflection argument \cite[Lemma~5.2]{Krisman2026chi2}, with prices $2fm$ at
the integer sites and $(2k+1)f$ at the poles; the numerator $\binom{u-1}A^2$ and the excess of
poles do not change the leading term.

\begin{remark}\label{rem:worse-sign}
Summing \cref{prop:p-divides-Y} over $p\mid q^2-4\sigma$ and \cref{prop:p-divides-q} over $p\mid
q$ gives the terms $-\frac{(2-d)^2}4\log(q^2-4\sigma)$ and $\frac13\log q$ of
$\aconst_\sigma(q)$. As $q^2-4<q^2+4$, the constant $\aconst_-(q)$ (the case of $\Ti_2$) is the
smaller one, and it suffices to compare the analytic exponent with $\aconst_-(q)$.
\end{remark}

\section{Ordinary primes}\label{sec:ordinary}

This section proves the lower bound $\Iprof_{\mathcal C}h^2$ for the contribution of the
ordinary primes $p\nmid2q(q^2-4\sigma)$ to $\sum_p\vG p(T)\log p$ (\cref{thm:ordinary}); it is the
largest term of $\aconst_\sigma(q)$ and is used in \cref{thm:arithmetic}. Primes $p<\delta h$ are
handled by the crude bound of \cref{prop:uniform} and primes $p>2K$ cost nothing
(\cref{lem:large-primes}). The work is in the range $p\asymp h$, where we bound $\vG p(T)$ by a
small optimisation problem over the residue classes modulo $p$ (\cref{subsec:two-bounds}), pass to
a limiting profile in $p/h$ (\cref{subsec:profile}), and integrate it against $\log p$.

Throughout, $p$ is ordinary, so $w$ and $1-w$ are $p$-adic units.

\begin{lemma}\label{lem:large-primes}
If $p>2K$ is ordinary, then $\vG p(T)\ge0$.
\end{lemma}

\begin{proof}
All factorials in $\nrm$ involve integers smaller than $p$, so $v_p(\nrm)=0$. The polynomial
$D_K$ is monic with $p$-integral coefficients, so the quotient of $u^{i+j+1}N_A(u)$ by $D_K$ has
$p$-integral coefficients and so has its derivative; by \cref{lem:moment-sums}(c) its moment sum
is $p$-integral. The residues $x_k^{i+j+1}N_A(x_k)/D_K'(x_k)$ are $p$-integral because
$D_K'(x_k)=\pm k!(K-1-k)!$, and $\tau_k$ involves only the odd numbers $2j+1<2K<p$.
\end{proof}

The interesting range is $p\asymp h$. Fix $\delta>0$. For $h$ large, every prime
$p\ge\delta h$ satisfies $p^2>3h>2K$ (both constructions have $2K<3h$), and we assume this from
now on. Then every difference
$l-k$ of pole indices and every odd number below $2K$ has $p$-adic valuation at most one.

\subsection{Residue classes and local functionals}\label{subsec:distribution}

\emph{Heuristic.} Modulo $p$ the poles fall into residue classes. In the coordinate $u=a+pz$
around a class $a$, the poles of that class are at distance $O(1)$ while all others are units,
so the weight $\varphi$ has a definite $p$-adic size on each class. Splitting the series over
the residue classes of $m$ (\cref{lem:distribution}) turns $\mu_X$ into a sum of local
functionals, one per class, which are integral (\cref{lem:local-integrality}).


For $a\in\{1,\dots,p\}$ let $\Pi_a=\{k<K:x_k\equiv a\pmod p\}$, $c_a=\abs{\Pi_a}$, and for
$k\in\Pi_a$ let $z_k=(x_k-a)/p\in\Z_p$; these are pairwise incongruent modulo $p$, because
$\abs{l-k}<K<p^2$. Let $\rest{\Z_p}$ be the ring of restricted power series
$\sum_jg_jz^j$ with $g_j\in\Z_p$, $g_j\to0$, and $\rest{\Q_p}=\rest{\Z_p}[1/p]$. Put $v=w^p$;
then $v\equiv w\pmod p$, so $v$ and $1-v$ are units. Let $M_r(v)=\sum_{n\ge0}n^rv^n\in\Z_p$ be
the rational moment ($M_0=1/(1-v)$, and $M_r=\sum_{i=1}^r S(r,i)\,i!\,v^i/(1-v)^{i+1}$ with
Stirling numbers $S(r,i)$), and define the continuous linear map
\[
  \mathsf M_v\colon\rest{\Q_p}\to\Q_p,\qquad \mathsf M_v\Bigl(\sum_jg_jz^j\Bigr)=\sum_jg_jM_j(v),
\]
which maps $\rest{\Z_p}$ into $\Z_p$. Let $V_a=\rest{\Q_p}\oplus\bigoplus_{k\in\Pi_a}\Q_p\cdot
(z-z_k)^{-1}$ and define $\ell_a\colon V_a\to\Q_p[X]$ by $\ell_a(G)=\mathsf M_v(G')$ for
$G\in\rest{\Q_p}$ and, for $k\in\Pi_a$,
\begin{equation}\label{eq:ell-pole}
  \ell_a\Bigl(\frac1{z-z_k}\Bigr)=\frac{p^2}{w^a}\Bigl[\mu_X\Bigl(\frac1{u-x_k}\Bigr)
  -\sum_{\substack{b=1\\b\ne a}}^{p}\frac{w^b}{p}\,
  \mathsf M_v\Bigl(\frac{\mathrm d}{\mathrm dz}\,\frac{1}{b+pz-x_k}\Bigr)\Bigr].
\end{equation}
For $R\in\Rsp$ let $R_a(z)=R(a+pz)$. If $R$ is a polynomial, $R_a\in\Q_p[z]$; if
$R=1/(u-x_k)$ then $R_a=p^{-1}(z-z_k)^{-1}$ for $k\in\Pi_a$, and otherwise $a-x_k$ is a unit and
$R_a\in\rest{\Z_p}$. Thus $R_a\in V_a$.

\begin{lemma}[distribution identity]\label{lem:distribution}
For every $R\in\Rsp$, $\displaystyle\mu_X(R)=\sum_{a=1}^{p}\frac{w^a}{p}\,\ell_a(R_a)$.
\end{lemma}

\begin{proof}
For a polynomial $g$ one has $\sum_{m\ge1}g(m)w^m=\sum_{a=1}^pw^a\sum_{n\ge0}g(a+pn)w^{pn}$ in
$\Q[[w]]$; both sides are rational functions of $w$, so the identity holds at our $w$, in
$\Q_p$. Apply it to $g=R'$ for a polynomial $R$, using $(R_a)'(z)=pR'(a+pz)$. For
$R=1/(u-x_k)$ with $k\in\Pi_{a_0}$ the identity is equivalent to the definition~\eqref{eq:ell-pole}
of $\ell_{a_0}(1/(z-z_k))$.
\end{proof}

\begin{lemma}[integrality]\label{lem:local-integrality}
Let $a\in\{1,\dots,p\}$ and $P_a(z)=\prod_{k\in\Pi_a}(z-z_k)$.
\begin{enumerate}[(a)]
\item For $k\in\Pi_a$, $\ell_a(1/(z-z_k))\in\Z_p[X]$.
\item If $N\in\rest{\Z_p}$, then $G=N/P_a\in V_a$ and $\ell_a(G)\in\Z_p[X]$.
\end{enumerate}
\end{lemma}

\begin{proof}
(a) In $p^2\mu_X(1/(u-x_k))=-4\sigma^ky^{2k+1}(p^2X-p^2\tau_k)$ the number $p^2\tau_k$ is
$p$-integral, since $v_p(2j+1)\le1$ for $2j+1<2K$. In the sum over $b\ne a$,
$p^{-1}\frac{\mathrm d}{\mathrm dz}(b+pz-x_k)^{-1}=-(b+pz-x_k)^{-2}$ lies in $\rest{\Z_p}$.
(b) Division of $z^j$ by the monic polynomial $P_a\in\Z_p[z]$ has integral quotient and
remainder, so $N=QP_a+N_0$ with $Q\in\rest{\Z_p}$ and $N_0\in\Z_p[z]$ of degree less than $c_a$.
Then $G=Q+\sum_{k\in\Pi_a}\frac{N_0(z_k)}{P_a'(z_k)}\frac1{z-z_k}$, and $P_a'(z_k)$ is a unit
because the $z_k$ are pairwise incongruent. Apply~(a).
\end{proof}

\emph{Local weights.} Let $Z_a=\{1\le j\le A:j\equiv a\pmod p\}$ and $n_a=\abs{Z_a}$. For $k\notin\Pi_a$ the factor
$a+pz-x_k$ is a unit of $\rest{\Z_p}$, and for $j\notin Z_a$ so is $a+pz-j$. Hence
\begin{equation}\label{eq:local-weight}
  \varphi(a+pz)=p^{\,2n_a-c_a}\,\frac{N_a(z)}{P_a(z)},\qquad
  N_a(z)=(a+pz)\prod_{j\in Z_a}\Bigl(z-\frac{j-a}p\Bigr)^2\,E_a(z),
\end{equation}
with $E_a\in\rest{\Z_p}^\times$, so $N_a\in\rest{\Z_p}$. We call
$\omega_a=2n_a-c_a-1$ the \emph{base cost} of the class $a$.

\subsection{Two lower bounds}\label{subsec:two-bounds}

\begin{proposition}[jet bound]\label{prop:jet}
$\vG p(\Delta)\ge\mathcal J_p$, where $\mathcal J_p$ is the minimum of
$\sum_a\sum_{j<t_a}(\omega_a+2j)$ over all $(t_a)_{a}$ with $t_a\ge0$ and $\sum_at_a=h$.
\end{proposition}

\begin{proof}
For polynomials $f$, $g$ of degree less than $h$ write $f(a+pz)=\sum_{i<h}f^{(a)}_ip^iz^i$; for
$f=u^r$ the Taylor coefficients $f^{(a)}_i=\binom ri a^{r-i}$ are integers. By
\cref{lem:distribution} and~\eqref{eq:local-weight},
\[
  \mu_X(\varphi u^ru^s)=\sum_{a=1}^p\sum_{i,j<h}(u^r)^{(a)}_i(u^s)^{(a)}_j\,\Lambda^{(a)}_{ij},
  \qquad
  \Lambda^{(a)}_{ij}=w^ap^{\,i+j-1+2n_a-c_a}\,\ell_a\Bigl(\frac{N_a(z)z^{i+j}}{P_a(z)}\Bigr),
\]
and $\vG p(\Lambda^{(a)}_{ij})\ge i+j+\omega_a$ by \cref{lem:local-integrality}(b). So the matrix of
$\Delta$ is $\mathbf J^{\mathsf T}\Lambda\mathbf J$ with $\mathbf J$ integral (rows indexed by
pairs $(a,i)$) and $\Lambda=\diag(\Lambda^{(a)})$ block diagonal. In~\eqref{eq:cauchy-binet} a
nonzero term selects $t_a$ rows and $t_a$ columns in each block, with distinct orders $i$ and
distinct orders $j$; each term of the Leibniz expansion of such a minor has valuation at least
$t_a\omega_a+2\binom{t_a}2=\sum_{j<t_a}(\omega_a+2j)$.
\end{proof}

Let $\epsilon_a=1$ if $\Pi_a$ contains an index $k$ with $2k+1<p$, and $\epsilon_a=0$
otherwise. Two elements of $\Pi_a$ differ by a multiple of $p$, so at most one index of $\Pi_a$
is below $(p-1)/2$, and it is the smallest one.

\begin{proposition}[pole bound]\label{prop:raw}
$\vG p(\Delta)\ge\mathcal R_p$, where $\mathcal R_p$ is the minimum of
$\sum_a\bigl[\sum_{j<t_a}(\omega_a+2j)+2\epsilon_a\one_{t_a=c_a}\bigr]$ over all $(t_a)$ with
$0\le t_a\le c_a$ and $\sum_at_a\le h$.
\end{proposition}

\begin{proof}
Write $\varphi u^s=Q_s+\sum_kx_k^s\rho_k/(u-x_k)$, $\rho_k=x_kN_A(x_k)/D_K'(x_k)$. The quotient
$Q_s$ of $u^{s+1}N_A$ by the monic $D_K\in\Z_p[u]$ is integral, so the matrix $\mathbf P=
[\mu_X(Q_{i+j})]$ is integral, and the matrix of $\Delta$ is
$\mathbf P+\sum_k\gamma_k\vvec(x_k)\vvec(x_k)^{\mathsf T}$ with
$\gamma_k=\rho_k\mu_X(1/(u-x_k))$. For $k\in\Pi_a$: $v_p(N_A(x_k))\ge2n_a$; exactly $c_a-1$
factors $l-k$ of $D_K'(x_k)$ are divisible by $p$, each exactly once; and $\mu_X(1/(u-x_k))$ has
Gauss valuation at least $-2$, and at least $0$ if $2k+1<p$. Hence
$\vG p(\gamma_k)\ge\omega_a+2\cdot\one_{2k+1<p}$.

Write the matrix as $[I_h\mid V]\diag(\mathbf P,\Gamma)[I_h\mid V]^{\mathsf T}$ with
$V=[\vvec(x_k)]_k$ and $\Gamma=\diag(\gamma_k)$. A nonzero term of~\eqref{eq:cauchy-binet}
selects the same set $\mathcal K$ of poles on both sides, $\abs{\mathcal K}\le h$, and some columns of
$I_h$. Its middle factor has valuation at least $\sum_{k\in\mathcal K}\vG p(\gamma_k)$. Each outer
factor is, up to sign, a minor of $V$ with columns $\mathcal K$; as a polynomial in the nodes it is
alternating with integer coefficients, hence divisible by $\prod_{k<l\in\mathcal K}(x_l-x_k)$, and
the $t_a=\abs{\mathcal K\cap\Pi_a}$ nodes of each class contribute valuation $\binom{t_a}{2}$. If
$t_a=c_a$, the index with $2k+1<p$, when it exists, is selected.
\end{proof}

The normalising factor contributes, by Legendre's formula and $p^2>3h>2K>h$,
\begin{equation}\label{eq:vp-S}
  v_p(\nrm)=h\Bigl(\floor{\tfrac Kp}-2\floor{\tfrac Ap}\Bigr)-2\sum_{i<h}\floor{\tfrac ip}
  =h\Bigl(\floor{\tfrac Kp}-2\floor{\tfrac Ap}-2Q_p+\tfrac phQ_p(Q_p+1)\Bigr),
\end{equation}
with $Q_p=\floor{h/p}$. Hence, for ordinary $p\ge\delta h$ and $h$ large,
\begin{equation}\label{eq:vp-T-combined}
  \vG p(T)\ \ge\ \max\bigl(-2h,\ v_p(\nrm)+\max(\mathcal J_p,\mathcal R_p)\bigr).
\end{equation}

\subsection{The profile}\label{subsec:profile}

\emph{Heuristic.} For $p=xh$ the residue classes come in a bounded number of types, according to
the numbers of poles and zeros they contain and to whether their first pole is improved, and the
number of classes of each type is $h$ times the length of an interval. The minima
$\mathcal J_p$ and $\mathcal R_p$ are then $h$ times the values of two filling problems whose
capacities are these lengths. We use the filling problems in their dual form, which makes both the
comparison with $\mathcal J_p$, $\mathcal R_p$ and the convexity in $x$ immediate.

\begin{lemma}[filling problems]\label{lem:filling}
Let $(e_i)_{i\in I}$ be real costs with capacities $m_i\ge0$, finitely many below any given
bound, and let $M>0$. For every selection $0\le y_i\le m_i$ with $\sum_iy_i=M$ and every
$\lambda\in\R$,
\[
  \sum_iy_ie_i\ \ge\ \lambda M-\sum_im_i\pos{\lambda-e_i},
\]
and the same holds for $\sum_iy_i\le M$ if $\lambda\le0$. Assume $\sum_im_i\ge M$. Equality holds for the
greedy selection of mass $M$ (cheapest costs first) with $\lambda$ its largest cost; in the second
case, for the greedy selection restricted to the negative costs (of mass $\min(M,\sum_{e_i<0}m_i)$)
with $\lambda$ the minimum of its largest cost and $0$.
\end{lemma}

\begin{proof}
$\sum_iy_ie_i=\lambda\sum_iy_i-\sum_iy_i(\lambda-e_i)$ and
$y_i(\lambda-e_i)\le m_i\pos{\lambda-e_i}$; if $\sum y_i\le M$ and $\lambda\le0$, then
$\lambda\sum_iy_i\ge\lambda M$. For the greedy selection of mass $M$ every $i$ with $e_i<\lambda$ is used
fully and every $i$ with $e_i>\lambda$ not at all, so both inequalities are equalities. In the
second case, if the negative costs have total capacity at least $M$, the greedy selection of mass
$M$ uses only negative costs and $\lambda\le0$ is its largest cost, so the first case applies;
otherwise take $y_i=m_i$ for $e_i<0$ and $y_i=0$ otherwise, and $\lambda=0$: both sides equal
$\sum_{e_i<0}m_ie_i$.
\end{proof}

\emph{The model.} Let $\kappa=K/h$ and $\alpha=A/h$ be fixed, and let $x>0$. Put
$c_0=\floor{\kappa/x}$, $r_\kappa=\kappa-c_0x$, $n_0=\floor{\alpha/x}$, $r_\alpha=\alpha-n_0x$, and
for $t\in[0,x)$
\begin{gather*}
  c(t)=c_0+\one_{t<r_\kappa},\qquad n(t)=n_0+\one_{((x/2-t)\bmod x)<r_\alpha},\\
  \epsilon(t)=\one_{t<x/2}\,\one_{c(t)>0},\qquad \omega(t)=2n(t)-c(t)-1.
\end{gather*}
Here $t$ stands for $k_0/h$, where $k_0$ is the least pole index in a class (see the proof of
\cref{prop:profile}). For $\lambda\in\R$ put
\[
  F_\lambda(\omega)=\sum_{j\ge0}\pos{\lambda-\omega-2j},\qquad
  \widetilde F_\lambda(\omega,c,\epsilon)=\sum_{0\le j<c}\pos{\lambda-\omega-2j-2\epsilon\one_{j=c-1}},
\]
and define
\begin{align*}
  J(x)&=\sup_{\lambda\in\R}\Bigl[\lambda-\int_0^xF_\lambda(\omega(t))\diff t\Bigr],\\
  R(x)&=\sup_{\lambda\le0}\Bigl[\lambda-\int_0^x\widetilde F_\lambda\bigl(\omega(t),c(t),
  \epsilon(t)\bigr)\diff t\Bigr],\\
  N(x)&=\floor{\kappa/x}-2\floor{\alpha/x}-2\floor{1/x}+x\floor{1/x}\bigl(\floor{1/x}+1\bigr),\\
  B(x)&=\max\bigl(-2,\ N(x)+\max(J(x),R(x))\bigr).
\end{align*}
The slots of the model take finitely many cost values $\omega(t)+2j$ below any bound, since
$\omega$ takes finitely many values; grouping the slots by cost value turns the continuous problems
into problems of the form of \cref{lem:filling}, with capacities the total length of $t$ with the
given cost. By \cref{lem:filling}, $J(x)$ is the least cost of filling mass $1$ with capacity $\diff t$ at each
cost $\omega(t)+2j$ (the \emph{jet slots}), and $R(x)$ the least cost of filling mass at most $1$
with capacity $\diff t$ at each cost $\omega(t)+2j+2\epsilon(t)\one_{j=c(t)-1}$, $j<c(t)$ (the
\emph{pole slots}); this primal form is the one evaluated in Appendix~\ref{app:profile-certificate}.

\begin{proposition}\label{prop:profile}
For every $\delta>0$ there is $C_\delta$ such that for all large $h$ in the construction and all
ordinary primes $\delta h\le p\le2K$, $\vG p(T)\ge h\,B(p/h)-C_\delta$.
\end{proposition}

\begin{proof}
Let $x=p/h\in[\delta,2\kappa]$. By~\eqref{eq:vp-S}, $v_p(\nrm)=hN(x)$ exactly, so
by~\eqref{eq:vp-T-combined} it suffices to show that the minima $\mathcal J_p$ and $\mathcal R_p$
of \cref{prop:jet,prop:raw} satisfy $\mathcal J_p\ge hJ(x)-C_\delta$ and
$\mathcal R_p\ge hR(x)-C_\delta$.

\emph{Exact class data.} For $0\le k_0<p$ let $x_{k_0}=-k_0-\frac12$ (a pole only if $k_0<K$).
These $p$ points lie in distinct classes modulo $p$, so the classes are indexed by $k_0$; the class
of $k_0$ contains the poles $x_k$ with $k\equiv k_0\pmod p$, $k<K$. Since $K/p=\kappa/x$ and
$K-p\floor{K/p}=hr_\kappa$, their number is $c(k_0/h)$ (for $K<p\le2K$ one has $c_0=0$,
$r_\kappa=\kappa$ and the classes with $k_0\ge K$ contain no pole, in accordance with
$c(t)=0$ for $t\ge r_\kappa$). Its first pole has
$2k_0+1<p$ if and only if $k_0\le(p-3)/2$, which agrees with $\one_{k_0/h<x/2}$ except for
$k_0=(p-1)/2$. Its zeros are the $j\in[1,A]$ with $j\equiv x_{k_0}\equiv\frac{p-1}2-k_0$; if
$\varrho\in\{1,\dots,p\}$ denotes this residue, their number is $n_0+\one_{\varrho\le hr_\alpha}$,
while $h\bigl((x/2-k_0/h)\bmod x\bigr)$ equals $\varrho+\frac12$ for $\varrho<p$ and $\frac12$ for
$\varrho=p$. Since $\varrho$ and $hr_\alpha$ are integers, the two indicators agree unless
$\varrho=hr_\alpha$ or $\varrho=p$. Altogether, at most three classes have data
$(c_a,n_a,\epsilon_a)$ different from $(c,n,\epsilon)(k_0/h)$.

\emph{Counting.} The model data are constant on each of the at most four intervals (open, closed or
half-open) into which the three interior points $x/2$, $r_\kappa$ and $(x/2-r_\alpha)\bmod x$ cut
$[0,x)$ (the set where $n(t)=n_0+1$ is $(s,x/2]$ cyclically, closed on the right), and an interval
of length $v-u$ contains between $h(v-u)-1$ and $h(v-u)+1$ of the points $k_0/h$; this gives an error of at most
$1$ per piece. Hence, for every $\lambda$,
\[
  \sum_aF_\lambda(\omega_a)\le h\int_0^xF_\lambda(\omega(t))\diff t+8\max_\omega F_\lambda(\omega),
\]
the maximum over $\omega\in[-c_0-2,\,2n_0+1]$, the range of all base costs; similarly for
$\widetilde F_\lambda$.

\emph{Comparison.} $\mathcal J_p$ is the least sum of $h$ elements of the multiset
$\{\omega_a+2j\}$, so \cref{lem:filling} gives $\mathcal J_p\ge h\lambda-\sum_aF_\lambda(\omega_a)$
for every $\lambda$. The supremum defining $J(x)$ is attained at the largest cost $\lambda^*$ used
by the greedy filling; as the levels $j<\lceil1/x\rceil$ have total capacity at least $1$,
$\lambda^*\le2n_0+1+2\lceil1/x\rceil$. With $\lambda=\lambda^*$ we obtain $\mathcal J_p\ge
hJ(x)-8\max_\omega F_{\lambda^*}(\omega)$, and $F_{\lambda^*}(\omega)\le(\lambda^*+c_0+4)^2/4$ is
bounded in terms of $\delta$, because $c_0\le\kappa/\delta$, $n_0\le\alpha/\delta$ and
$1/x\le1/\delta$. In the same way $\mathcal R_p$ is the least sum of at most $h$ elements of the
multiset of pole slots, and \cref{lem:filling} applies with $\lambda\le0$. The supremum over
$\lambda\le0$ defining $R(x)$ is attained, since the bracket is concave and piecewise affine in
$\lambda$ with slope $1$ for $\lambda<\min\omega-2$ (all terms vanish there), so a maximiser
$\lambda^*\le0$ exists. For $\lambda\le0$ and $\omega\ge-c_0-2$ each of the at most $c_0+1$
terms of $\widetilde F_\lambda$ is at most $c_0+2$, hence
$\widetilde F_\lambda\le(c_0+1)(c_0+2)$, and $\mathcal R_p\ge hR(x)-8(c_0+1)(c_0+2)$.
\end{proof}

\begin{lemma}[piecewise convexity]\label{lem:convexity}
Let $\mathfrak B$ be the set of points $1/j$, $\kappa/j$, $2\kappa/j$, $\alpha/j$, $2\alpha/j$,
$2(\kappa+\alpha)/j$, $j\ge1$ (only the first three if $\alpha=0$). On every open interval of
$(0,\infty)\setminus\mathfrak B$ the function $B$ is convex. Consequently
$\int_a^bB\ge(b-a)B(\frac{a+b}2)$ for consecutive points $a<b$ of $\mathfrak B\cup\{\delta\}$.
\end{lemma}

\begin{proof}
Let $U$ be an open interval disjoint from $\mathfrak B$. The integers $\floor{1/x}$, $c_0$, $n_0$
change only at the points $1/j$, $\kappa/j$, $\alpha/j$, so they are constant on $U$; hence $N$ is
affine on $U$, and $x/2$, $x$ and $r_\kappa=\kappa-c_0x$ are affine. The point
$s=(x/2-r_\alpha)\bmod x$ equals $x/2-r_\alpha$ or $x/2-r_\alpha+x$ according as
$r_\alpha\le x/2$ or not, and $r_\alpha=x/2$ means $x=2\alpha/(2n_0+1)\in\mathfrak B$, so $s$ is
affine on $U$ too. Two of the cut points $0$, $x/2$, $r_\kappa$, $s$, $x$ coincide only if
$r_\kappa=0$ ($x=\kappa/c_0$), $r_\kappa=x/2$ ($x=2\kappa/(2c_0+1)$), $s=0$ ($r_\alpha=x/2$),
$s=x/2$ ($r_\alpha=0$, $x=\alpha/n_0$) or $s=r_\kappa$ (then $\kappa+\alpha\in x/2+x\Z$, so
$x=2(\kappa+\alpha)/(2m+1)$); the remaining coincidences ($r_\kappa=x$, $s=x$, $x/2=0$) are
impossible. All these values lie in $\mathfrak B$, so on $U$ the cut points keep their order, the
model data $(c,n,\epsilon)$ are constant on each piece, and the lengths of the pieces are affine
in $x$. For fixed $\lambda$ the functions $x\mapsto\lambda-\int_0^xF_\lambda(\omega(t))\diff t$
and $x\mapsto\lambda-\int_0^x\widetilde F_\lambda\diff t$ are therefore affine on $U$, so $J$ and $R$, as
suprema of affine functions, are convex. They are finite: $R\le0$ and
$J\le2n_0+1+2\lceil1/x\rceil$. Hence $B$ is convex and bounded on $U$. A convex function on $(a,b)$ lies above
its supporting line at $\frac{a+b}2$, which gives the integral inequality.
\end{proof}

\subsection{The ordinary contribution}\label{subsec:ordinary-contribution}

\begin{theorem}\label{thm:ordinary}
For every $\delta>0$,
\[
  \liminf_{h\to\infty}\frac1{h^2}\sum_{p\ \mathrm{ordinary}}\vG p(T)\log p\ \ge\
  \int_\delta^{2\kappa}B(x)\diff x-2\delta .
\]
With $\delta=1/2000$ the right-hand side is at least $\Iprof_\Gen=-2.113164137$ for
$(\kappa,\alpha)=(7/6,0)$ and at least $\Iprof_\Band=-1.987089582$ for $(49/40,1/20)$.
\end{theorem}

\begin{proof}
Split the ordinary primes into $p<\delta h$, $\delta h\le p\le2K$ and $p>2K$. The last range
contributes a nonnegative amount by \cref{lem:large-primes}. For the first, by
\cref{prop:uniform},
\[
  \sum_{p<\delta h}2hL_p\log p\le2h\,\theta(\delta h)+2h\sqrt{3h}\log(3h)=2\delta h^2+o(h^2),
\]
because $L_p=1$ for $\sqrt{3h}<p<\delta h$, with $\theta(x)=\sum_{p\le x}\log p\sim x$. In the middle
range, \cref{prop:profile} and $\sum_{p\le2K}\log p=O(h)$ reduce the claim to
$\sum_{\delta h\le p\le2K}B(p/h)\log p\ge h\int_\delta^{2\kappa}B-o(h)$. The function $B$ is bounded
on $[\delta,2\kappa]$ and continuous outside the finite set $\mathfrak B\cap[\delta,2\kappa]$,
hence Riemann integrable; approximating it from below by step functions and using
$\theta(\beta h)-\theta(\beta'h)=(\beta-\beta')h+o(h)$ for fixed $\beta'<\beta$ (the prime number
theorem) gives the claim. The finitely many excluded primes do not affect the limit. The numerical
values are lower bounds for the midpoint sums of \cref{lem:convexity} over the
$5999$ (respectively $11798$) intervals cut out by $\mathfrak B$ in $[\delta,2\kappa]$, computed
exactly in rational arithmetic (Appendix~\ref{app:profile-certificate}).
\end{proof}

\begin{remark}\label{rem:profile-k1}
For $(\kappa,\alpha)=(1,0)$ the same computation gives $\int_\delta^2B-2\delta=-1.7169166\ldots
=\frac12-\frac{133}{60}-\frac\delta2$. Since $\log\nrm=\frac12h^2+O(h\log h)$ when $\kappa=1$,
this is Krisman's generic constant $\frac{133}{60}$ \cite[Thm.~4.5]{Krisman2026chi2} up to the
cut-off: at $\kappa=1$ the profile reproduces his class bound, and the gain for $\Gen$ does not
come from a sharper estimate at ordinary primes. What the profile adds is that it is defined for
every $\kappa$ and $\alpha$.
\end{remark}

\section{The size of \texorpdfstring{$T(\xi)$}{T(ξ)}}\label{sec:decay}

This section proves the upper bound $\limsup h^{-2}\log\abs{T(\xi)}\le\bconst$
(\cref{thm:decay}), the second input of the criterion; it is combined with the arithmetic bound in
\cref{sec:proofs}. \emph{Heuristic.} By~\eqref{eq:real-sites}, $T(\xi)$ is a sum over
configurations of $h$ points on the summation sites, at most two per site, each configuration
weighted by a product of Vandermonde factors and of the weights $w^m\varphi(m)$. This is a discrete
logarithmic gas in an external field, and its size is controlled by comparing the configurations
with a fixed step density $\rho$; the comparison costs only $O(h\log h)$ thanks to a discrete
energy inequality. The method is Krisman's \cite[\S6--7]{Krisman2026chi2}; we redo it because the
external field now contains the terms coming from $K\ne h$ and from the zero band.

Throughout, $K=\kappa h$ and $A=\alpha h$ with $\kappa\ge1$, $0\le\alpha<1$ and
$1\le d=\kappa-2\alpha<2$ fixed.

\subsection{Statement}\label{subsec:decay-statement}

A \emph{comparison density} is a function $\rho\colon\R\to[0,2]$, constant on the cells of a
finite partition of $[\alpha,1]$ into intervals, vanishing outside $[\alpha,1]$, with
$\int\rho=1$. Put
\begin{align*}
  U(x)&=\int\rho(t)\log\abs{x-t}\diff t,\qquad E(\rho)=\int U(x)\rho(x)\diff x,\\
  V(x)&=x\log W+3\psi(x)-\psi(x+\kappa)-2\psi(x-\alpha)\qquad(x\ge\alpha),\\
  G(x)&=2U(x)+V(x),\qquad C=\kappa\log\kappa-2\alpha\log\alpha+\tfrac32-d,
\end{align*}
with $0\log0=0$. Note that $3\psi(x)-\psi(x+\kappa)-2\psi(x-\alpha)=-\int_x^{x+\kappa}\log t\diff
t+2\int_{x-\alpha}^x\log t\diff t$: the first term comes from the poles and the second from the
zero band.

\begin{theorem}\label{thm:decay}
Let $\rho$ be a comparison density and $\lambda\in\R$, and suppose that there is $\eta>0$ with
$G(x)-\lambda\le-\eta x$ for all $x\ge1$. Then
\[
  \limsup_{h\to\infty}\frac{1}{h^2}\log\abs{T(\xi)}\ \le\ \bconst(\rho,\lambda):=
  C-E(\rho)+\lambda+2\int_\alpha^1\pos{G(x)-\lambda}\diff x .
\]
\end{theorem}

\begin{lemma}[tail condition]\label{lem:tail}
For $x\ge1$, $G(x)\le x\log W+(2-d)\log x-\frac\kappa2\log\bigl(1+\frac\kappa{2x}\bigr)$.
Consequently the hypothesis of \cref{thm:decay} holds if
\[
  \tau_0:=\log W-\tfrac\kappa2\log\bigl(1+\tfrac\kappa2\bigr)-\lambda<0\quad\text{and}\quad
  \tau_1:=1+\log W+\tfrac{\kappa^2}4<0 .
\]
\end{lemma}

\begin{proof}
For $x\ge1$ and $t\in[\alpha,1]$, $\log\abs{x-t}\le\log x$, so $U(x)\le\log x$. Next,
$\int_{x-\alpha}^x\log t\diff t\le\alpha\log x$, and by concavity of $\log$,
$\int_x^{x+\kappa}\log t\diff t\ge\frac\kappa2(\log x+\log(x+\kappa))\ge\kappa\log
x+\frac\kappa2\log(1+\frac\kappa{2x})$ (the middle term is
$\kappa\log x+\frac\kappa2\log(1+\frac\kappa x)$; we keep the weaker form, which is all that is used).
This gives the first inequality. Since $0<2-d\le1$,
$(2-d)\log x\le x-1$; the function $g(x)=-\frac\kappa2\log(1+\frac\kappa{2x})$ has
$g'(x)=\frac{\kappa^2}{4x^2+2\kappa x}\le\frac{\kappa^2}4$ for $x\ge1$. Hence
$G(x)-\lambda\le\tau_0+(x-1)\tau_1\le-\eta x$ with $\eta=\min(-\tau_0,-\tau_1)$.
\end{proof}

\subsection{Three lemmas and the proof}\label{subsec:decay-lemmas}

\begin{lemma}[minors with derivative columns]\label{lem:derivative-minors}
Let $h\ge2$ and let $s_1,\dots,s_h$ be integers, each value occurring at most twice. Let
$\mathcal E$ be the $h\times h$ matrix whose $i$-th column is $\vvec(s_i)$ or $\vvec'(s_i)$, where
for a value occurring twice one column of each kind is taken. Then
\[
  \abs{\det\mathcal E}\le h^h\Bigl(1+\frac2h\Bigr)^{h^2}\prod_{i<j,\ s_i\ne s_j}\abs{s_i-s_j}.
\]
\end{lemma}

\begin{proof}
Write each derivative column as $\vvec'(s)=\frac1{2\pi i}\oint_{\abs{\zeta-s}=1/h}
\vvec(\zeta)(\zeta-s)^{-2}\diff\zeta$. By multilinearity $\det\mathcal E$ is an integral of
Vandermonde determinants $\prod_{i<j}(\zeta_j-\zeta_i)$, with $\zeta_i=s_i$ for value columns
and $\abs{\zeta_i-s_i}=1/h$ otherwise, against a kernel of total mass at most $h$ per derivative
column. If $s_i\ne s_j$ then $\abs{\zeta_i-\zeta_j}\le\abs{s_i-s_j}+2/h\le(1+2/h)\abs{s_i-s_j}$;
if $s_i=s_j$ then $\abs{\zeta_i-\zeta_j}=1/h\le1$.
\end{proof}

\begin{lemma}[a discrete energy inequality]\label{lem:polya}
For every finitely supported real sequence $(\nu_s)_{s\in\Z}$ with $\sum_s\nu_s=0$,
\[
  \sum_{s\ne t}\nu_s\nu_t\log\abs{s-t}\ \le\ \log2\sum_s\nu_s^2 .
\]
\end{lemma}

\begin{proof}
We may assume the support lies in $\{0,\dots,M-1\}$ with $M\ge2$. Define $f$ on $\Z_{\ge0}$ by
$f(0)=\log(2M)$, $f(d)=\log(M/d)$ for $1\le d\le M$ and $f(d)=0$ for $d\ge M$. Its second
differences $f(d-1)-2f(d)+f(d+1)$ are $0$ at $d=1$, $\log\frac{d^2}{d^2-1}$ for $2\le d<M$,
$\log\frac M{M-1}$ at $d=M$, and $0$ beyond, so $f$ is convex and decreases to $0$, whence
$f(d)=\sum_{k\ge1}\bigl(f(k-1)-2f(k)+f(k+1)\bigr)\pos{k-d}$. The kernel
$\pos{k-\abs{s-t}}$ counts the windows of $k$ consecutive integers containing $s$ and $t$, so it is
positive semidefinite, and therefore $\sum_{s,t}\nu_s\nu_tf(\abs{s-t})\ge0$. Expanding,
this sum equals $\log2\sum_s\nu_s^2+\log M\,(\sum_s\nu_s)^2-\sum_{s\ne t}\nu_s\nu_t\log\abs{s-t}$.
\end{proof}

This is the Pólya-type inequality of \cite[Lemma~6.4]{Krisman2026chi2}.

\begin{lemma}[discretisation]\label{lem:discretisation}
Let $\rho$ be a comparison density, $m_s=h\int_{(s-1)/h}^{s/h}\rho$ for $s\in\Z$, and
\[
  P_m(s)=\sum_{t\ne s}m_t\log\abs{s-t},\qquad E_m=\sum_{s\ne t}m_sm_t\log\abs{s-t}.
\]
Then $0\le m_s\le2$, $\sum_sm_s=h$, $m_s=0$ unless $A<s\le h$, and for every integer $s$,
\[
  P_m(s)\le hU(s/h)+h\log h+6+2\log h,\qquad E_m\ge h^2E(\rho)+h^2\log h-4h(1+\log h).
\]
\end{lemma}

\begin{proof}
Let $\varrho(u)=\rho(u/h)$, which has mass $m_t$ on $I_t=[t-1,t]$, and note
$hU(s/h)=\int\varrho(u)\log\abs{s-u}\diff u-h\log h$ and
$h^2E(\rho)+h^2\log h=\iint\varrho(u)\varrho(u')\log\abs{u-u'}\diff u\diff u'$. For the first
inequality compare, cell by cell, $m_t\log\abs{s-t}$ with $\int_{I_t}\varrho(u)\log\abs{s-u}\diff
u$. For $t<s$ the discrete term is the smaller one. For $t=s$ the discrete term is absent and
the integral is at least $-2\int_0^1\abs{\log v}\diff v=-2$. For $t=s+e$, $e\ge1$, the excess is at
most $2\int_{e-1}^{e}(\log e-\log v)\diff v$, which is $2$ for $e=1$ and at most $2/(e-1)$ for
$e\ge2$; summing over $e\le h$ gives at most $4+2\log h$. For the second inequality, cells with
$t=t'$ contribute a nonpositive amount to the double integral, and for $\abs{t-t'}=e\ge1$ we
have $\log\abs{u-u'}\le\log(e+1)\le\log e+1/e$, so the double integral is at most
$E_m+\sum_tm_t\sum_{1\le e\le h}2\cdot2/e$.
\end{proof}

\begin{proof}[Proof of \cref{thm:decay}]
By~\eqref{eq:real-sites}, $T(\xi)=\nrm\lim_{M\to\infty}\det H^{(M)}$, where $H^{(M)}$ is the sum
over $A<m\le M$. Apply~\eqref{eq:cauchy-binet}. A nonzero term selects, at each site $s$, the
same number $n_s\in\{0,1,2\}$ of slots on both sides, with $\sum_sn_s=h$; each multiplicity
vector $n=(n_s)$ arises from at most $4^h$ pairs of slot sets. For $s>A$,
$\varphi'/\varphi(s)=1/s+\sum_{j\le A}2/(s-j)-\sum_k1/(s-x_k)$ has absolute value at most
$1+2A+K\le3h$, so every minor of the middle factor is bounded by $\prod_sb(s)^{n_s}$ with
$b(s)=3hW^s\abs{\varphi(s)}$. With \cref{lem:derivative-minors},
\begin{equation}\label{eq:decay-step1}
  \abs{\det H^{(M)}}\le4^hh^{2h}e^{4h}\sum_n\prod_sb(s)^{n_s}\prod_{s<t}\abs{s-t}^{2n_sn_t}.
\end{equation}

\emph{Energy.} Let $m$ be as in \cref{lem:discretisation} and $\nu=n-m$. Then $\sum_s\nu_s=0$ and
$\sum_s\nu_s^2\le2\sum_s(n_s^2+m_s^2)\le8h$. Expanding $n=m+\nu$ and using
$\sum_{s\ne t}\nu_sm_t\log\abs{s-t}=\sum_sn_sP_m(s)-E_m$, \cref{lem:polya} and
\cref{lem:discretisation} give
\begin{align*}
  \sum_{s\ne t}n_sn_t\log\abs{s-t}
  &=-E_m+2\sum_sn_sP_m(s)+\sum_{s\ne t}\nu_s\nu_t\log\abs{s-t}\\
  &\le h^2\log h-h^2E(\rho)+2h\sum_sn_sU(s/h)+O(h\log h).
\end{align*}

\emph{Weights.} For $s>A$ the function $t\mapsto\log(s-t)$ decreases on $[0,A]$, so
$\sum_{j=1}^A\log(s-j)\le\int_0^A\log(s-t)\diff t=\psi(s)-\psi(s-A)$; and by concavity
$\log D_K(s)=\sum_{k<K}\log(s+k+\frac12)\ge\int_0^K\log(s+t)\diff t=\psi(s+K)-\psi(s)$. Since
$\psi(hx)=h\psi(x)+hx\log h$, with $x=s/h$,
\[
  \log b(s)\le hV(x)-(K-2A)\log h+\log s+\log(3h).
\]

\emph{Summation.} Combining, the logarithm of the summand in~\eqref{eq:decay-step1} is at most
\[
  (1-d)h^2\log h-h^2E(\rho)+h^2\lambda+\sum_sn_s\bigl(h(G(s/h)-\lambda)+\log s\bigr)+O(h\log h).
\]
For $s\le h$ we use $n_sh(G-\lambda)\le2h\pos{G-\lambda}$ and $\log s\le\log h$; for $s>h$ the
tail hypothesis gives $h(G(s/h)-\lambda)\le-\eta s$. Since every multiplicity vector comes from at
least one ordered $h$-tuple of sites,
\[
  \sum_n\prod_s\bigl(se^{h(G(s/h)-\lambda)}\bigr)^{n_s}\le
  \exp\Bigl(2h\sum_{A<s\le h}\pos{G(s/h)-\lambda}\Bigr)\Bigl(h^2+\sum_{s>h}se^{-\eta s}\Bigr)^h,
\]
uniformly in $M$. Finally Stirling's formula gives
$\log\nrm=(d-1)h^2\log h+Ch^2+O(h\log h)$, the terms $h^2\log h$ cancel, and since $G$ is
continuous on $[\alpha,1]$ the Riemann sums $h^{-1}\sum_{A<s\le h}\pos{G(s/h)-\lambda}$ converge
to $\int_\alpha^1\pos{G-\lambda}$.
\end{proof}

\subsection{Larger denominators}\label{subsec:shift}

The exponent depends on $q$ only through $\log W=\log4-2\log q$, and only in the term $x\log W$
of $V$.

\begin{proposition}[shift]\label{prop:shift}
Let $\alpha=0$, let $q_0\le q$ be odd, and let $(\rho,\lambda)$ satisfy the conditions of
\cref{lem:tail} for $W_0=4/q_0^2$. Put $\tau=2\log(q/q_0)$. Then $(\rho,\lambda-\tau/2)$
satisfies them for $W=4/q^2$, and
$\bconst_q(\rho,\lambda-\tau/2)\le\bconst_{q_0}(\rho,\lambda)-\frac12\log(q/q_0)$.
\end{proposition}

\begin{proof}
We have $G_q(x)=G_{q_0}(x)-\tau x$, so $G_q-(\lambda-\tau/2)=G_{q_0}-\lambda-\tau(x-\frac12)$ and
$\pos{G_q-\lambda+\tau/2}\le\pos{G_{q_0}-\lambda}+\tau\pos{\frac12-x}$. Integrating over $[0,1]$
costs at most $2\cdot\tau/8$, while $\lambda$ decreases by $\tau/2$. The quantities $\tau_0$
and $\tau_1$ of \cref{lem:tail} decrease by $\tau/2$ and $\tau$.
\end{proof}

This is Krisman's shift lemma \cite[Lemma~7.1]{Krisman2026chi2}. The bound of
\cref{thm:decay} depends on $\abs w$ only, so it is the same for $\chi_2$ and $\Ti_2$; Krisman
observes \cite[Remark~6.7]{Krisman2026chi2} that for $w<0$ the true determinants are smaller,
which we do not use.

\section{Nonvanishing at prime indices}\label{sec:nonvanishing}

This section proves that $T_n(r)\ne0$ for every rational $r$ and every large prime $n$
(\cref{thm:nonvanishing} with \cref{prop:local-dets-general,prop:local-dets-band}), the third input
of the criterion, used in \cref{sec:proofs}. \emph{Heuristic.} When $h$ and $K$ are fixed multiples
of a prime $p$, every residue class modulo $p$ contains the same numbers of poles and zeros. In a
basis of polynomials adapted to the classes, the distribution identity of \cref{sec:ordinary}
makes the scaled matrix block triangular modulo $p$, and each diagonal block reduces to one of
three fixed rational matrices that do not depend on $p$. If their determinants are nonzero,
$\Delta(r)$ is a $p$-adic unit times an explicit power of $p$.

Let $t\ge1$, $\gamma\ge1$ and $a_0\ge0$ be fixed integers, and for a prime $p$ consider the
determinant with
\[
  h=tp,\qquad K=\gamma p,\qquad A=a_0p .
\]
The general construction is $(t,\gamma,a_0)=(6,7,0)$ with $n=p$, and the band construction is
$(40,49,2)$. Put $e_0=\gamma-2a_0$ and $Q=\sigma y^2=\sigma q^2/4$, so that $w=1/Q$.

\subsection{The local matrices and the statement}\label{subsec:local-matrices}

For $s\in\{0,1\}$ let $\mathcal V_s$ be the $\Q$-vector space spanned by $\Q[z]$ and the
fractions $1/(z+l+\frac12)$, $s\le l<s+\gamma$, and define $\ell_Q\colon\mathcal V_s\to\Q$ by
\begin{equation}\label{eq:ellQ}
  \ell_Q(g)=\sum_{r\ge0}[z^r]g'\cdot M_r(1/Q)\quad(g\in\Q[z]),\qquad
  \ell_Q\Bigl(\frac1{z+l+\frac12}\Bigr)=4\sum_{i<l}\frac{Q^{\,l-i}}{(2i+1)^2},
\end{equation}
where $M_r(v)=\sum_{n\ge0}n^rv^n$ is the rational function of \cref{subsec:distribution}. For
$(s,e)\in\{(0,0),(1,0),(1,1)\}$ let $L_{s,e}(Q)$ be the $t\times t$ matrix with entries
\[
  \ell_Q\Bigl(\frac{z^{i+j}(z+1)^e\prod_{i'<a_0}(z-i')^2}{\prod_{l=s}^{s+\gamma-1}(z+l+\frac12)}
  \Bigr),\qquad 0\le i,j<t .
\]
Its entries are rational numbers depending on $q$, $\sigma$ and $(t,\gamma,a_0)$ only.

\begin{theorem}\label{thm:nonvanishing}
Assume that the three determinants $\det L_{s,e}(Q)$ are nonzero. Then for every $r\in\Q$ and
every sufficiently large prime $p$ one has $\Delta(r)\ne0$, hence $T(r)\ne0$. More precisely,
$v_p(\Delta(r))=-pt(e_0+2-t)+t$.
\end{theorem}

For $\Gen$ the valuation is $-18p+6$. We checked this prediction against exact computations of
$\Delta(0)$ for $q=83$ at $p=17$ ($\sigma=-1$, $h=102$) and $p=19$ ($\sigma=1$, $h=114$), where
$v_p(\Delta(0))=-300$ and $-336$ respectively (Appendix~\ref{app:consistency}).

\begin{proof}[Proof of \cref{thm:nonvanishing}]
Let $p$ be a prime not dividing $2q(q^2-4\sigma)$ or the denominator of $r$, with $p>2\gamma+1$
(so $p^2>2K$), and not dividing the numerators and denominators of the three determinants. Then $p\nmid
Q(Q-1)$, and the entries of the $L_{s,e}(Q)$, which are built from the numbers $(2i+1)^{-2}$,
$i<\gamma$, differences of poles, powers of $Q$ and the moments $M_r(1/Q)$ with denominators
powers of $Q-1$, are $p$-integral because $p>2\gamma+1$ and $p\nmid Q(Q-1)$. All but finitely
many primes qualify. We use the
notation of \cref{sec:ordinary}; every residue class $a\in\{1,\dots,p\}$ contains exactly
$\gamma$ poles and $a_0$ zeros of $N_A$, namely $a,a+p,\dots,a+(a_0-1)p$.

\emph{Step 1: a basis adapted to the classes.} The ideals $((u-a)^t)$, $1\le a\le p$, of
$\Z_p[u]$ are pairwise comaximal, because $a-b$ is a unit for $a\ne b$. By the Chinese remainder
theorem there are $f_{a,i}\in\Z_p[u]$ of degree less than $tp$, for $1\le a\le p$ and $0\le
i<t$, with $f_{a,i}\equiv(u-a)^i\pmod{(u-a)^t}$ and $f_{a,i}\equiv0\pmod{(u-b)^t}$ for $b\ne a$.
They form a $\Z_p$-basis of the polynomials of degree less than $h$ with $\Z_p$ coefficients, so
their coefficient matrix $C$ has $\det C\in\Z_p^\times$, and by \cref{lem:change-of-basis} the
Gram matrix $\mathbf G=[\mu_r(\varphi f_{a,i}f_{b,j})]$ has $\det\mathbf G=(\det C)^2\Delta(r)$.
Write $f_{a,i}(c+pz)=p^iF_{a,i,c}(z)$; then $F_{a,i,c}\in\Z_p[z]$, $F_{a,i,a}\equiv z^i$ and,
for $c\ne a$, $F_{a,i,c}\in p^{t-i}\Z_p[z]\subset p\Z_p[z]$.

\emph{Step 2: local weights.} By~\eqref{eq:local-weight}, for $c\ne p$,
$\varphi(c+pz)=p^{-e_0}\Psi_c(z)$ with $\Psi_c=N_c/P_c$, $N_c\in\rest{\Z_p}$ and
\[
  N_c\equiv\lambda_c\prod_{i'<a_0}(z-i')^2\pmod p,\qquad
  P_c(z)=\prod_{l\in J_c}\bigl(z+l+\tfrac12\bigr),
\]
where, for $c\ne p$, $\lambda_c$ is the unit $c\prod_{k\notin\Pi_c}(c-x_k)^{-1}\prod_{j\notin Z_c}(c-j)^2$
(for $c=p$ the factor $c$ is omitted: $\lambda_p=\prod_{k\notin\Pi_p}(p-x_k)^{-1}\prod_{j\notin Z_p}(p-j)^2$).
Here $J_c=\{0,\dots,\gamma-1\}$ if $c\le\frac{p-1}2$ and $J_c=\{1,\dots,\gamma\}$ if
$c\ge\frac{p+1}2$: indeed $x_k-c=-(2k+2c+1)/2$, and $2k+2c+1=(2l+1)p$ for $0\le k<\gamma p$ gives
exactly these ranges of $l$, with $z_k=-l-\frac12$ \emph{exactly}. For $c=p$ the numerator
$u=p(1+z)$ supplies one more factor $p$: $\varphi(p+pz)=p^{1-e_0}\Psi_p(z)$ with
$N_p\equiv\lambda_p(1+z)\prod_{i'<a_0}(z-i')^2$.

\emph{Step 3: the scaled matrix.} By \cref{lem:distribution},
$\mathbf G_{(a,i),(b,j)}=\sum_cw^cp^{\,i+j-1-e_0}\ell_c(\tilde\varphi_cF_{a,i,c}F_{b,j,c})$, where
$\tilde\varphi_c=\Psi_c$ for $c\ne p$ and $\tilde\varphi_p=p\Psi_p$. Let $\mathbf M$ be obtained from
$\mathbf G$ by multiplying row $(a,i)$ by $p^{e_0+1-i}$, and moreover by $p^{-1}$ if $a=p$, and
column $(b,j)$ by $p^{-j}$. Then
\[
  \mathbf M_{(a,i),(b,j)}=p^{-\one_{a=p}}\sum_{c=1}^pw^c\,\ell_c\bigl(\tilde\varphi_cF_{a,i,c}F_{b,j,c}\bigr),
\]
which is $p$-integral by \cref{lem:local-integrality}(b): for $a\ne p$ directly; for $a=p$ because
$F_{p,i,c}\in p\Z_p[z]$ when $c\ne p$, and $\tilde\varphi_p\in p\,V_p$.

\emph{Step 4: block triangularity.} Modulo $p$, $F_{a,i,c}\equiv\one_{a=c}z^i$, and $\ell_c$ is
$\Z_p$-linear on the lattice of \cref{lem:local-integrality}(b), so a term with $c\ne a$ or $c\ne b$
vanishes modulo $p$ unless it is divided by $p$. Hence for $a\ne p$ the row $(a,i)$ of
$\mathbf M$ is, modulo $p$, supported in the columns $(a,j)$, with entries
$w^a\ell_a(\Psi_az^{i+j})$. For $a=p$, the entries in columns $(p,j)$ receive from $c\ne p$ terms
in $p^{2t-i-j-1}\Z_p\subset p\Z_p$, so they are $w^p\ell_p(\Psi_pz^{i+j})$ modulo $p$. Ordering the
classes with $p$ last, $\mathbf M$ is block lower triangular modulo $p$, and
\begin{equation}\label{eq:det-product}
  \det\mathbf M\equiv\prod_{a=1}^{p}\det\bigl[w^a\ell_a(\Psi_az^{i+j})\bigr]_{i,j<t}\pmod p .
\end{equation}

\emph{Step 5: reduction of the local functionals.} We claim that for $N\in\rest{\Z_p}$,
$\ell_a(N/P_a)\equiv\ell_Q(\bar N/P_a)\pmod p$, where $\bar N$ is the polynomial obtained by
reducing the coefficients of $N$ and $\ell_Q$ is evaluated on a lift: choose
$\tilde N\in\Z_{(p)}[z]$ with reduction $\bar N$; then $\ell_Q(\tilde N/P_a)\in\Z_{(p)}$ because
$p>2\gamma+1$ and $p\nmid Q(Q-1)$, and its class modulo $p$ does not depend on the lift, since
changing $\tilde N$ by $pg$ changes the value by $p\,\ell_Q(g/P_a)\in p\Z_{(p)}$. For the restricted-series
part this follows from $v=w^p\equiv w=1/Q\pmod p$ and the fact that $M_r$ has denominator
$(1-v)^{r+1}$, a unit. For a pole $k\in\Pi_a$ with $z_k=-l-\frac12$, the sum over $b\ne a$
in~\eqref{eq:ell-pole} is divisible by $p^2$, and $p^2r\equiv0$, so
$\ell_a(1/(z-z_k))\equiv4\sigma^ky^{2k+1}w^{-a}p^2\tau_k$. In $p^2\tau_k$ only the terms with
$p\mid2i'+1$ survive; writing $2i'+1=(2i+1)p$, they are those with $i<l$ (because
$2k+1=(2l+1)p-2a$ and $1\le a\le p$), and they contribute
$\sigma^{ip+(p-1)/2}y^{-(2i+1)p}/(2i+1)^2$. Using $\sigma^p=\sigma$, $y^p\equiv y$, $w^{-a}=\sigma^a
y^{2a}$, $y^{2k+2a+1}=y^{(2l+1)p}$ and $k+a+\frac{p-1}2=(l+1)p-1\equiv l\pmod2$, we obtain
\[
  \ell_a\Bigl(\frac{1}{z-z_k}\Bigr)\equiv4\sum_{i<l}\sigma^{l+i}y^{2l-2i}(2i+1)^{-2}
  =4\sum_{i<l}\frac{Q^{\,l-i}}{(2i+1)^2}\pmod p,
\]
which is~\eqref{eq:ellQ}. Since the decomposition of $N/P_a$ into a restricted series and
principal parts has integral coefficients (\cref{lem:local-integrality}(b)) and commutes with
reduction, the claim follows.

\emph{Step 6: conclusion.} By Steps 2 and 5 the $a$-th block in~\eqref{eq:det-product} is
congruent to $w^a\lambda_a$ times $L_{s_a,e_a}(Q)$, with $s_a=\one_{a>(p-1)/2}$ and $e_a=\one_{a=p}$.
Hence $\det\mathbf M$ is a $p$-adic unit. As $\det\mathbf M=p^{\,pt(e_0+1)-pt(t-1)-t}\det\mathbf
G$ and $\det C$ is a unit, $v_p(\Delta(r))=-pt(e_0+2-t)+t$.
\end{proof}

Krisman's Lemma~P \cite[Thm.~8.1]{Krisman2026chi2} follows the same pattern with $h=3\ell$,
$K=h$, the Lagrange basis at the poles and three $3\times3$ blocks. With $K>h$ there is no
Lagrange basis at the poles, which is why we use the Taylor jets of Step~1.

\subsection{The local determinants}\label{subsec:local-dets}

\begin{proposition}\label{prop:local-dets-general}
For $(t,\gamma,a_0)=(6,7,0)$ and $(s,e)=(0,0),(1,0),(1,1)$,
\[
  (Q-1)^{24}\det L_{s,e}(Q)=c_{s,e}\,Q^{\,a_{s,e}}(Q-1)^{b_{s,e}}\,\Pi_{s,e}(Q)
\]
with nonzero rationals $c_{s,e}$, $(a_{s,e},b_{s,e})=(9,20),(11,20),(13,18)$, and primitive
integer polynomials $\Pi_{s,e}$ of degrees $16,20,20$ whose roots all have modulus less than
$189$. In particular $\det L_{s,e}(\pm q^2/4)\ne0$ for every odd $q\ge29$.
\end{proposition}

\begin{proof}
The entries of $(Q-1)^4L_{s,e}(Q)$ are polynomials in $Q$: the derivative of the polynomial part
has degree at most $3$, so only $M_0,\dots,M_3$ occur, and the pole values are polynomials. The
determinants were computed exactly, twice and independently (Appendix~\ref{app:local-certificates}): once
symbolically in $Q$, and once by exact interpolation from $90$ rational values, with agreement.
The bound on the roots is Cauchy's bound $1+\max_i\abs{\pi_i/\pi_{\deg}}$, evaluated exactly; it
is less than $23.19$, $136.73$ and $188.09$ for the three polynomials. (Numerically all roots have modulus less
than $1.88$; this is not used.) The statement for every odd $q\ge29$ uses only the Cauchy bound.
\end{proof}

\begin{proposition}\label{prop:local-dets-band}
For $(t,\gamma,a_0)=(40,49,2)$, $q=81$ and $\sigma=\pm1$, the six rational numbers
$\det L_{s,e}(\pm81^2/4)$ are nonzero.
\end{proposition}

\begin{proof}
They were computed exactly in rational arithmetic (Appendix~\ref{app:local-certificates}); their
numerators have between $7233$ and $7548$ decimal digits. As an independent check, their
reductions modulo $191$ are $171$, $138$, $106$ for $\sigma=1$ and $47$, $24$, $90$ for
$\sigma=-1$, in agreement with a separate computation carried out entirely modulo $191$.
\end{proof}

\section{Proof of the main theorem}\label{sec:proofs}

This section combines the three estimates through the criterion of \cref{prop:criterion}: it
collects the local bounds into the arithmetic constant (\cref{thm:arithmetic}), compares it with
the certified size exponent for each construction (\cref{thm:general,thm:band}), and records the
status of every step of the proof (\cref{subsec:step-status}).

\begin{theorem}[arithmetic]\label{thm:arithmetic}
Let $\mathcal C\in\{\Gen,\Band\}$, $q\ge5$ odd and $\sigma=\pm1$. Along the $n$ with $T_n\ne0$,
$\liminf h^{-2}\sum_p\vG p(T_n)\log p\ge\aconst_\sigma(q)$, with $\aconst_\sigma(q)$ as in
\cref{thm:three-estimates}.
\end{theorem}

\begin{proof}
Sum \cref{thm:prime-two} (prime $2$), \cref{prop:p-divides-q} over the primes dividing $q$,
\cref{prop:p-divides-Y} over the primes dividing $q^2-4\sigma$, and \cref{thm:ordinary}; these sets
of primes are disjoint and the first three are finite. Note $\sum_{p\mid q}v_p(q)\log p=\log q$.
\end{proof}


\begin{theorem}\label{thm:general}
For every odd $q\ge83$ and $\sigma=\pm1$, the number $\xi$ of~\eqref{eq:basic-notation} is
irrational.
\end{theorem}

\begin{proof}
Let $(\rho_{83},\lambda_{83})$ be the comparison data of Appendix~\ref{app:decay-certificates}. They
satisfy the conditions of \cref{lem:tail} for $W=4/83^2$, and
\[
  \bconst_{83}:=\bconst(\rho_{83},\lambda_{83})\le-3.161708,\qquad
  \aconst_-(83)\ge-3.154846 .
\]
For odd $q\ge83$ put $\lambda_q=\lambda_{83}-\log(q/83)$. By \cref{prop:shift},
$\bconst(\rho_{83},\lambda_q)\le\bconst_{83}-\frac12\log(q/83)$ for $W=4/q^2$. On the arithmetic
side, \cref{rem:worse-sign} and $(q^2+4)/(83^2+4)\le(q/83)^2$ give
\[
  \aconst_\sigma(q)\ge\aconst_-(q)\ge\aconst_-(83)+\Bigl(\frac13-2\cdot\frac{25}{144}\Bigr)
  \log\frac q{83}=\aconst_-(83)-\frac1{72}\log\frac q{83}.
\]
Hence $\bconst(\rho_{83},\lambda_q)-\aconst_\sigma(q)\le-0.006862-\frac{35}{72}\log(q/83)<0$. By
\cref{thm:arithmetic,thm:decay} the first two hypotheses of \cref{prop:criterion} hold for the
family $\Gen$. By \cref{thm:nonvanishing,prop:local-dets-general}, for every $r\in\Q$ we have
$T_p(r)\ne0$ for all large primes $p$; in particular $T_n\ne0$ for infinitely many $n$.
\Cref{prop:criterion} gives $\xi\notin\Q$.
\end{proof}

\begin{theorem}\label{thm:band}
For $q=81$ and $\sigma=\pm1$, the number $\xi$ is irrational.
\end{theorem}

\begin{proof}
With the comparison data $(\rho_{81},\lambda_{81})$ of Appendix~\ref{app:decay-certificates}, which
satisfy \cref{lem:tail} for $W=4/81^2$, we have $\bconst\le-3.374176$ and
$\aconst_-(81)\ge-3.368498$, so $\bconst<\aconst_\sigma(81)$ for both signs. Nonvanishing follows
from \cref{thm:nonvanishing,prop:local-dets-band}, and \cref{prop:criterion} concludes.
\end{proof}

\begin{proof}[Proof of \cref{thm:main}]
Combine \cref{thm:general,thm:band} with $\sigma=1$ ($\chi_2$) and $\sigma=-1$ ($\Ti_2$), and use
that both functions are odd.
\end{proof}

\subsection{Status of the individual steps}\label{subsec:step-status}

We use three labels. \emph{Proved}: an elementary argument, given in full, of a kind we regard
as routine. \emph{Candidate}: an argument given here in full (or, where indicated, in outline)
that has not yet been checked by anyone independent of the AI systems involved.
\emph{Certified (AI-internal)}: a finite computation whose result was reproduced by a second
implementation written without reference to the first, both by AI systems; this is not an
independent verification. \Cref{thm:main} as a whole is a candidate.
\begin{center}
\small
\begin{tabular}{@{}p{0.44\textwidth}p{0.48\textwidth}@{}}
\toprule
Step & Status\\
\midrule
Moment sums, pole series, change of basis, criterion
(\cref{lem:moment-sums,lem:pole-series,lem:change-of-basis,prop:criterion}) & proved; the final
deduction is formalised in Lean, conditionally on the three estimates\\
Prime $2$ (\cref{thm:prime-two}) & candidate (new); consistent with exact valuations for
$h\le72$\\
Primes dividing $q(q^2-4\sigma)$ (\cref{sec:exceptional}) & candidate (adapted from Krisman);
consistent with exact valuations at small $h$\\
Ordinary primes (\cref{thm:ordinary}) & candidate (new); finite bounds consistent with exact
valuations for $h\le42$\\
Profile integrals $\Iprof_\Gen$, $\Iprof_\Band$ & certified (AI-internal)\\
Size bound (\cref{thm:decay}) & candidate (adapted from Krisman)\\
Comparison data $\bconst_{83}$, $\bconst_{81}$ & certified (AI-internal)\\
Nonvanishing mechanism (\cref{thm:nonvanishing}) & candidate (new); predicted valuation confirmed
at $p=17,19$\\
Local determinants (\cref{prop:local-dets-general,prop:local-dets-band}) & certified
(AI-internal)\\
\bottomrule
\end{tabular}
\end{center}

\begin{remark}[points for scrutiny]\label{rem:scrutiny}
In our judgement the steps most in need of an independent check are: (i) the comparison of the
exact class data with the limiting profile in \cref{prop:profile} and the breakpoint argument of
\cref{lem:convexity}, on which the use of midpoint sums as lower bounds depends; (ii) the bookkeeping of derivative columns and
of the middle factor in the proof of \cref{thm:prime-two}, which is new and on which the constant
$-140/99$ rests; (iii) the reduction of the local functionals in Step~5 of the proof of
\cref{thm:nonvanishing}; (iv) the implementations of the certificates in
Appendix~\ref{app:certificates}, which were written by AI systems. The margins (about $0.0069$ and $0.0057$
per $h^2$) are small compared with the individual terms of \cref{tab:constants}, so an error of
a few thousandths in any constant would invalidate the corresponding case.
\end{remark}

\section{Limitations: the case \texorpdfstring{$q=79$}{q = 79}}\label{sec:limitations}

The argument stops at $q=81$. This section explains why, quantifies the obstruction at $q=79$ for
the band construction, and lists what would be needed to go further; nothing in it is used in the
proof. Its statements concern the specific estimates of this paper, not the determinants
themselves, nor the irrationality of $\chi_2(2/79)$ or $\Ti_2(2/79)$.

\emph{The general construction.} For $\Gen$ the arithmetic constant decreases only like
$-\frac1{72}\log q$ while the analytic exponent improves like $-\frac12\log q$, so the margin
shrinks as $q$ decreases. With the comparison density recorded for $q=81$ (mesh $1/240$, as in
Appendix~\ref{app:decay-certificates}) the ball-arithmetic verifier gives the upper bound
$\bconst-\aconst_-(81)\le+0.005332$ for this density. This is an upper bound for one density found
by a numerical optimiser, not a lower bound for the method: it shows that $\Gen$ does not reach
$q=81$ with these data, and we have not proved that no comparison density does better.

\emph{The band construction at $q=79$: a limitation of these bounds.} The numbers below are
computed in floating point or with a step density and are \emph{not certified}; the section
describes what the estimates of this paper can and cannot give, not a property of the
determinants. Weak duality bounds every exponent of \cref{thm:decay} from below: if $\mu$ is any
density with compact support in $[\alpha,\infty)$, $0\le\mu\le2$ and $\int\mu=1$, then
$E(\mu-\rho)\le0$ (the logarithmic energy of a compactly supported signed measure of total mass
zero and finite energy is nonpositive \cite[Lemma~I.1.8]{SaffTotik1997}; both $\mu$ and $\rho$
are bounded densities with compact support), so
\[
  C+E(\mu)+\int V\mu\ \le\ C-E(\rho)+\int G\mu\ \le\ \bconst(\rho,\lambda),
\]
using $\mu\le2$ and $G<\lambda$ on $[1,\infty)$. For the step density $\mu_{79}$ recorded in the
supplementary data, the left-hand side at $q=79$ is $-3.358974$ (evaluated in ball arithmetic,
which makes this number, but not the comparison below, rigorous).
On the arithmetic side we use the upper bound
$\int_0^{2\kappa}B\le-1.985157$ for $(\kappa,\alpha)=(49/40,1/20)$. It is obtained from
trapezoidal sums with the class types frozen at the midpoint ($B$ is convex on each interval of
\cref{lem:convexity}), evaluated in exact rational arithmetic, plus the contribution of
$(0,\delta)$ under the \emph{unproved} bound $B\le1$ on $(0,\delta)$, $\delta=1/2000$. We have not
proved $B\le1$ there (it should follow from an explicit estimate of the $O(1)$ quantities
involved); it is plausible but unproved. With this input the arithmetic
constant at $q=79$ is at most $-3.365088$ for $\sigma=1$ and $-3.365333$ for $\sigma=-1$ (floating
point). The resulting gap of more than $0.006$ \emph{indicates} that, in this geometry, no choice
of comparison density and no refinement of the numerical integration would make the estimates of
\cref{sec:prime-two,sec:exceptional,sec:ordinary,sec:decay} prove the case $q=79$; as long as
$B\le1$ on $(0,\delta)$ is unproved and the arithmetic inputs are floating-point values, this is
a heuristic obstruction for these bounds, not a theorem. The three numbers were recomputed
(Appendix~\ref{app:consistency}).

\emph{What would be needed.} Since the obstruction concerns sufficient estimates, it leaves
several possibilities open: other pole ratios or zero placements; lower bounds for the actual
valuations that are sharper than the slot profiles, for instance through cancellation between
residue classes; or, for $\Ti_2$, a size bound that sees the sign of $w$, which Krisman suggests
as the remaining lever for his own construction \cite[\S12]{Krisman2026chi2}. Exploratory
searches over positive half-integer zeros, reflected integer zeros and two-sided zero bands
found no favourable configuration at $q=79$, but they were not exhaustive and several stopped at
time limits; they are not evidence of impossibility.

\emph{The square construction.} With $K=h$ and the parity bound of \cref{sec:prime-two} in place
of Krisman's Lemma~11.1, the same method would close the comparison at $q=127$ but not at $q=125$,
(unverified) according to an earlier stage of our computations, which we have not re-verified.
This is not a literature threshold; we mention it only to indicate how much of the progress from
$175$ to $83$ is due to the prime $2$ and how much to the ratio $K/h$.


\appendix
\section{The finite computations}\label{app:certificates}

Three kinds of finite computations enter the proof. Each was first carried out by the code that
accompanies the candidate arguments, and then recomputed by a second implementation written
from the definitions in this paper (the programs listed in the section on reproducibility). Both implementations were written by
AI systems; agreement between them is what we call \emph{certified (AI-internal)}.

\subsection{The profile integrals}\label{app:profile-certificate}

For $(\kappa,\alpha)$ and $\delta=1/2000$, the breakpoints $\mathfrak B\cap[\delta,2\kappa]$ of
\cref{lem:convexity}, together with $\delta$ and $2\kappa$, cut $[\delta,2\kappa]$ into $5999$
intervals for $(7/6,0)$ and $11798$ intervals for $(49/40,1/20)$. At the midpoint $x$ of each
interval the class types, their capacities and $N(x)$ are rational, and $J(x)$, $R(x)$ are
values of finite linear programmes with rational data. The first implementation fills the
cheapest slots greedily after pre-filling the jet levels below $\floor{1/x}-2$ (respectively
$\floor{1/x}-4$), which is legitimate because every slot of those levels is cheaper than every
slot at level $\floor{1/x}-1$ or above, and the levels below $\floor{1/x}-1$ have total capacity
less than one. The second implementation does not pre-fill: it finds by bisection the least
integer cost $c^*$ such that the slots of cost at most $c^*$ have capacity at least one, using the
closed form $\sum_{\text{class}}\mathrm{cap}\cdot\#\{j:\omega+2j\le c\}$, and fills up to $c^*$.
All arithmetic is exact. Both give
\[
  \sum_{\text{intervals}}(b-a)\,B\bigl(\tfrac{a+b}2\bigr)-2\delta=
  \begin{cases}-2.11316413653\ldots&(\kappa,\alpha)=(7/6,0),\\
  -1.98708958191\ldots&(\kappa,\alpha)=(49/40,1/20),\\
  -1.71691666666\ldots&(\kappa,\alpha)=(1,0),\end{cases}
\]
in agreement to all digits; the stated constants $\Iprof_\Gen=-2113164137\cdot10^{-9}$ and
$\Iprof_\Band=-993544791\cdot(5\cdot10^{8})^{-1}$ are rational numbers below the exact sums. As a
measure of the slack in the midpoint rule, the trapezoidal upper bound of \cref{sec:limitations}
for $(49/40,1/20)$ is $-1.985157$, which exceeds the midpoint sum by $0.0019$.

\subsection{The comparison densities}\label{app:decay-certificates}

The densities $\rho_{83}$ and $\rho_{81}$ are step functions on the mesh $1/240$ with cell masses
$m_i/10^7$, $m_i\in\Z_{\ge0}$, $\sum m_i=10^7$, stored in the supplementary files
\path{certificates/general-construction-comparison-density.json} and
\path{certificates/band-construction-comparison-density.json} together with
$\lambda_{83}=-7.1118866$ and $\lambda_{81}=-7.6105436$. The density $\rho_{83}$ equals the cap $2$
on $[0,0.4)$ and decreases to $0$ on $[0.4,155/240]$; $\rho_{81}$ vanishes on $[0,1/20)$, equals
$2$ on $[1/20,106/240)$ and decreases on $[106/240,1]$. They were found by numerical optimisation;
only their stored values matter.

The quantities in $\bconst(\rho,\lambda)$ are evaluated as follows.
\begin{itemize}
\item $E(\rho)$ is a finite sum over pairs of cells of
$\iint_{\text{cell}\times\text{cell}}\log\abs{x-y}=g(d+s)+g(d-s)-2g(d)$, with $s=1/240$, $d$ the
difference of the left endpoints, and $g(t)=\frac12t^2\log\abs t-\frac34t^2$.
\item $G$ is a finite sum of terms $c_j\psi(x-a_j)$ with rational $c_j$ and $a_j$ on the grid
$\frac1{24000}\Z$, plus $x\log W$. On each grid cell of length $\ell=1/24000$, $G$ is at most the
larger of its endpoint values plus $\sum_j\abs{c_j}\epsilon_j$, where $\epsilon_j$ bounds the
deviation of $\psi(x-a_j)$ from its chord: $\ell^2/(8\,\mathrm{dist})$ if $a_j$ is at distance
$\mathrm{dist}\ge\ell$ from the cell, and $\ell/e$ if $a_j$ is an endpoint (because
$\abs{t\log(t/\ell)}\le\ell/e$ on $[0,\ell]$).
\item The integral $2\int(G-\lambda)_+$ is bounded by $2\ell$ times the sum of the positive parts
of these cell bounds minus $\lambda$, and the tail conditions of \cref{lem:tail} are checked
directly: $(\tau_0,\tau_1)=(-0.6076,-6.111)$ for $q=83$ and $(-0.0847,-6.027)$ for $q=81$.
\end{itemize}
The first implementation uses fixed-point logarithms with denominator $2^{100}$ and explicit error
budgets; the second uses ball arithmetic (Arb, $160$ bits) for every logarithm. They give, with all decimals rounded in the safe direction,
\begin{center}
\begin{tabular}{@{}lllll@{}}
\toprule
 & $E(\rho)$ & $\bconst\le$ & $\aconst_-\ge$ & $\bconst-\aconst_-\le$\\
\midrule
$q=83$, $\Gen$ & $-2.121571081$ & $-3.161708963$ & $-3.154845833$ & $-0.006863130$\\
$q=81$, $\Band$ & $-2.113189373$ & $-3.374176117$ & $-3.368497052$ & $-0.005679065$\\
\bottomrule
\end{tabular}
\end{center}
(second implementation; the first gives upper bounds for $\bconst$ larger by less than
$10^{-8}$, because it uses the cruder constant $\ell/2$ in place of $\ell/e$). The arithmetic
constant $\aconst_-$ uses the rational lower bounds $\Iprof$ of Appendix~\ref{app:profile-certificate}. A
third, floating-point evaluation with a midpoint rule on $48000$ points gives
$\bconst\approx-3.16191$ and $-3.37435$, slightly below the certified upper bounds, as it should.

\subsection{The local determinants}\label{app:local-certificates}

For $\Gen$ the entries of $(Q-1)^4L_{s,e}(Q)$ are polynomials in $Q$ with rational
coefficients. The first implementation computes the three determinants symbolically. The second
evaluates $(Q-1)^{24}\det L_{s,e}(Q)$ exactly at the $90$ integers $Q=2,\dots,91$, interpolates,
checks the interpolant at three further rational points (among them $Q=-83^2/4$), removes the
factors $Q$ and $Q-1$, and compares with the first: the integer polynomials $\Pi_{s,e}$ agree up
to a rational factor, with the same exponents $(a_{s,e},b_{s,e})$. The leading-term domination
$\sum_{i\ge1}\abs{\pi_{\deg-i}/\pi_{\deg}}R^{-i}<1$ at $R=83^2/4$, which is a second, separate proof of
nonvanishing for $\abs Q\ge R$ (hence for every $q\ge83$; either this or the Cauchy bound of
\cref{prop:local-dets-general} suffices for the proof), holds with left-hand sides $0.0022$, $0.0054$ and $0.0043$.

For $\Band$ at $q=81$ the $40\times40$ matrices have rational entries computed from the definition
\eqref{eq:ellQ}: polynomial division by the $49$ local poles, exact residues, and the moments
$M_r(1/Q)$ for $r\le34$. Their determinants are nonzero rational numbers; their reductions
modulo $191$ (a prime not dividing any denominator involved, since all the odd numbers and pole
differences occurring are smaller than $191$ and $Q$, $Q-1$ are units) are as stated in
\cref{prop:local-dets-band}, and coincide with the values obtained by the first implementation,
which works modulo $191$ throughout.

\section{Consistency checks at small sizes}\label{app:consistency}

The checks below compare the statements of the paper with exact computations of $T$ for small
$h$. They cannot confirm asymptotic statements, whose error terms are of the same size as the
differences observed, but they would have exposed errors in the finite forms of the bounds.
The polynomials $T$ were computed exactly by an implementation of the definitions in
\cref{sec:pencils} (interpolation of $\det[\mu_X(\varphi u^{i+j})]$ at $h+1$ rational points).

\subsection{The pole series}
\Cref{lem:pole-series} was tested on several entries $\mu_\xi(\varphi u^s)$ with the series
$\sum_mw^m(\varphi u^s)'(m)$ summed directly, with exact rational derivatives. For $\Gen$ at
$q=83$ and $\Band$ at $q=81$, $n=1$, both signs, the relative discrepancy is below $10^{-1240}$ at
$1500$ decimal digits (an earlier audit at $400$ digits found $10^{-60}$ and $10^{-145}$). For
$\Band$ the two parts of the affine entries are of size up to $10^{156}$ and cancel almost
completely at $X=\xi$; a first test at $120$ digits reported a spurious mismatch for this reason.

\subsection{The prime 2}
\Cref{tab:two-adic} compares $\vG2(T)$ with the finite bound of \cref{thm:prime-two} for $\Gen$ at
$q=83$ ($\vG2(T)$ is the same for both signs at these sizes).

\begin{table}[ht]
\centering
\caption{Exact $\vG2(T)$ for $\Gen$, $q=83$, against the finite bound of \cref{thm:prime-two}
and the leading term $-\frac{140}{99}h^2$. Last line: $\Band$, $q=81$, against
$-\frac{1847}{1100}h^2$.}\label{tab:two-adic}
\begin{tabular}{@{}rrrrr@{}}
\toprule
$h$ & $\vG2(T)$ & finite bound & leading term & $\vG2(T)/h^2$\\
\midrule
6 & $-40$ & $-43$ & $-50.9$ & $-1.111$\\
12 & $-172$ & $-180$ & $-203.6$ & $-1.194$\\
18 & $-405$ & $-417$ & $-458.2$ & $-1.250$\\
24 & $-730$ & $-748$ & $-814.5$ & $-1.267$\\
30 & $-1159$ & $-1177$ & $-1272.7$ & $-1.288$\\
36 & $-1696$ & $-1722$ & $-1832.7$ & $-1.309$\\
42 & $-2328$ & $-2358$ & $-2494.5$ & $-1.320$\\
48 & $-3051$ & $-3083$ & $-3258.2$ & $-1.324$\\
54 & $-4048$ & $-4080$ & $-4123.6$ & $-1.388$\\
60 & $-4816$ & $-4854$ & $-5090.9$ & $-1.338$\\
66 & $-5917$ & $-5966$ & $-6160.0$ & $-1.358$\\
72 & $-6998$ & $-7054$ & $-7330.9$ & $-1.350$\\
\midrule
40 ($\Band$) & $-2441$ & $-2462$ & $-2686.5$ & $-1.526$\\
\bottomrule
\end{tabular}
\end{table}

The ratios are not monotone, and all lie above $-140/99\approx-1.4141$. The entry at $h=54$ is an
outlier ($-1.388$, against $-1.324$ and $-1.338$ at $h=48,60$); it was recomputed on
26 November 2026 and gives the same value $-4048$. The finite bound shows the same dip
($-4080$), so it is a feature of the bound and not an error in the exact value; $h=54$ is the
size with $K=63=111111_2$, the largest binary digit sum of the $K$ in the table, which may explain the dip
through the low-order terms of \cref{thm:prime-two}; we have not isolated the mechanism.

For Krisman's determinant ($K=h$, $A=0$, $q=83$, $\sigma=1$) \cref{tab:two-adic-square} gives the
exact values in two normalisations: the un-normalised $\Delta=\det[\mu_X(\varphi u^{i+j})]$ and
$T=S\Delta$ with the normaliser $S$ of this paper, against the finite bound of
\cref{thm:prime-two} (valid for $T$).

\begin{table}[ht]
\centering
\caption{Exact $2$-adic Gauss valuations for $K=h$, $A=0$, $q=83$, $\sigma=1$; the limit of the
bound is $-\frac{13}{11}h^2$, that is $-1.1818h^2$. The columns with $/h^2$ show minus the
valuation divided by $h^2$.}\label{tab:two-adic-square}
\begin{tabular}{@{}rrrrrrr@{}}
\toprule
$h$ & $\vG2(\Delta)$ & $\vG2(T)$ & bound for $T$ & $-\vG2(\Delta)/h^2$ & $-\vG2(T)/h^2$ & $-$bound$/h^2$\\
\midrule
8 & $-66$ & $-42$ & $-48$ & $1.0312$ & $0.6562$ & $0.7500$\\
16 & $-287$ & $-223$ & $-236$ & $1.1211$ & $0.8711$ & $0.9219$\\
24 & $-663$ & $-583$ & $-604$ & $1.1510$ & $1.0122$ & $1.0486$\\
32 & $-1188$ & $-1028$ & $-1053$ & $1.1602$ & $1.0039$ & $1.0283$\\
40 & $-1882$ & $-1722$ & $-1760$ & $1.1763$ & $1.0762$ & $1.1000$\\
48 & $-2700$ & $-2492$ & $-2538$ & $1.1719$ & $1.0816$ & $1.1016$\\
56 & $-3685$ & $-3485$ & $-3533$ & $1.1751$ & $1.1113$ & $1.1266$\\
64 & $-4816$ & $-4432$ & $-4484$ & $1.1758$ & $1.0820$ & $1.0947$\\
\bottomrule
\end{tabular}
\end{table}

\subsection{Odd primes}
For $\Gen$ at $q=83$ and $h=6,12,\dots,42$, both signs, and for $\Band$ at $q=81$ and $h=40$, both
signs, every prime $p\le\max(4h,2K+2)+q$ was tested against the finite form of the relevant
bound: \cref{prop:p-divides-Y} for $p\mid q^2-4\sigma$, \cref{lem:large-primes} for $p>2K$,
\eqref{eq:vp-T-combined} with the exact class counts for $p^2>2K$, and \cref{prop:uniform}
otherwise. No bound failed. Typical margins: for $h=42$ the ordinary primes $p=41$ and $p=43$
have $\vG p(T)=-25$ and $-19$ against bounds $-26$ and $-20$; for $q^2-4=6885=3^4\cdot5\cdot17$ the
primes $3,5,17$ have $\vG p(T)=-1360,-422,-367$ against $-1560,-558,-390$. For $p=83$ one has
$\vG{83}(T)=h^2/3$ exactly for $h\le30$ and then slightly less ($0.3326h^2$ at $h=36$, $0.3206h^2$
at $h=54$, $0.3219h^2$ at $h=60$), as allowed by the term $-O(h\log h)$ in
\cref{prop:p-divides-q}. For $\Band$, $\vG3(T)=2106$ at $h=40$, against the leading term
$\frac43h^2=2133.3$.

\subsection{Nonvanishing}
For $\Gen$ at $q=83$, \cref{thm:nonvanishing} predicts $v_p(\Delta(0))=-18p+6$ for every
admissible prime $p$ with $h=6p$. Exact computation of $\Delta(0)$ gives $v_{17}(\Delta(0))=-300$
for $\sigma=-1$, $h=102$ (note $17\mid83^2-4$, so $p=17$ is excluded for $\sigma=1$), and
$v_{19}(\Delta(0))=-336$ for $\sigma=1$, $h=114$, as predicted. These are checks of the whole
mechanism of \cref{sec:nonvanishing}, including the reduction of the local functionals.

\subsection{Size}
The audit also computed $h^{-2}\log\abs{T(\xi)}$ in ball arithmetic. For $\Gen$ at $q=83$ it is
$-3.130$, $-3.111$, $-3.107$, $-3.107$, $-3.110$, $-3.111$, $-3.115$, $-3.116$ for
$h=12,\dots,54$ ($\sigma=1$), above the asymptotic bound $\bconst_{83}\le-3.1617$ by amounts
compatible with an error $O(\log h/h)$; for $\Band$ at $h=40$ it is $-3.377$ ($\sigma=1$) and
$-3.379$ ($\sigma=-1$), against $\bconst\le-3.3742$.

\subsection{The case \texorpdfstring{$q=79$}{q = 79}}
In floating point we recomputed the three numbers of \cref{sec:limitations}: the value
$-3.358974$ of $C+E(\mu_{79})+\int V\mu_{79}$, the trapezoidal bound $-1.985157$ for the profile
integral, and the resulting upper bounds $-3.365088$ and $-3.365333$ for the arithmetic
constant at $q=79$.


% back matter starts on a new page
\clearpage
\afbacksection{Reproducibility and formal verification}
\sloppy

\emph{Status of the result.} \Cref{thm:main} is a candidate. No human expert independent of the
AI systems that produced it has reviewed the argument. Each step carries one of the labels of
\cref{subsec:step-status}: \emph{proved} (routine elementary arguments), \emph{candidate}
(argued in this paper, not yet independently verified) or \emph{certified (AI-internal)} (a finite
computation reproduced by a second implementation, both written by AI systems). None of this is
an independent verification.

\emph{Computations.} Paths are relative to the directory of this contribution (AF-0003) in the
archive \url{https://github.com/axiomforge-math/mathematics}. The programs are in
\path{logs/experiments/scripts/}, the step densities and integer polynomials in
\path{logs/experiments/certificates/}, the recorded outputs in \path{logs/experiments/outputs/},
and every command with its outcome in \path{logs/verification.md}. The programs need Python~3.11
or later with \texttt{python-flint}~0.9.0 and \texttt{mpmath}~1.3.0
(\path{logs/experiments/requirements.txt}); from \path{logs/experiments/scripts/} the command
\texttt{python3 run\_checks.py all} runs all of them (the largest sizes take about half an hour
each). The programs of the form \path{verify_*.py} and \path{check_*.py} are the second
implementations, written from the definitions in this paper; none of them imports the code that
produced the candidate certificates. The data files are:
\begin{itemize}
\item \path{certificates/general-construction-comparison-density.json},
\item \path{certificates/band-construction-comparison-density.json},
\item \path{certificates/general-construction-local-polynomials.json}.
\end{itemize}
The programs and what they check:
\begin{itemize}
\item \path{verify_profile_integrals.py}: the profile integrals of \cref{thm:ordinary} and
\cref{rem:profile-k1} (exact rational arithmetic);
\item \path{verify_comparison_bounds.py}: the rigorous bounds for $\bconst$ and $\aconst_-$ of
Appendix~\ref{app:decay-certificates} (Arb ball arithmetic), including the tail conditions of
\cref{lem:tail};
\item \path{verify_local_determinants.py} and \path{compute_local_blocks.py}: the local
determinants of \cref{prop:local-dets-general,prop:local-dets-band};
\item the consistency checks of Appendix~\ref{app:consistency}: \path{check_finite_bounds.py},
\path{check_prime_index_valuation.py}, \path{check_series_identity.py},
\path{audit_small_sizes.py}, and \path{compute_two_adic_valuations.py} (the values of $\vG2(T)$
for $h\le72$ in \cref{tab:two-adic}); the exact pencils are implemented in
\path{define_half_integer_pencils.py}, the values for $K=h$, $A=0$ of
\cref{tab:two-adic-square} are obtained with the same implementation for $(h,K,A)=(h,h,0)$;
\item for \cref{sec:limitations}: \path{bound_profile_integral.py} and
\path{estimate_comparison_bounds.py} (numbers that are not certified).
\end{itemize}
The programs that first produced the certificates were also re-executed during the internal audit
and reproduce their logs; they are not part of the package. All of this was carried out by AI
systems.

\emph{Formal verification.} The Lean~4 package in \path{lean/} (Lake package
\texttt{axiomforge-af0003}, root module \texttt{AxiomForge.AF0003}, namespace
\texttt{AxiomForge.LegendreChi}) is pinned by \path{lean/lean-toolchain}
(\texttt{leanprover/lean4:v4.34.0-rc1}) and \path{lean/lakefile.toml} (Mathlib at the exact commit
recorded there). The \emph{main theorems are not formalised}: they are proved only conditionally,
as \path{AxiomForge.LegendreChi.irrational_legendreChi_two_div_of_estimates} and
\path{AxiomForge.LegendreChi.irrational_inverseTangentIntegral_two_div_of_estimates} (and
their versions for $-2/q$), which derive the irrationality for odd $q\ge81$ from a hypothesis that
packages the arithmetic estimate, the size estimate and the nonvanishing at prime indices, with
rounded rates. That hypothesis is exactly the unproved content of
\cref{sec:prime-two,sec:exceptional,sec:ordinary,sec:decay,sec:nonvanishing}; the Lean theorems
are therefore not formalisations of \cref{thm:main}. Unconditionally the package proves the
oddness of $\chi_2$ and $\Ti_2$ and the transfer between $x$ and $-x$, the irrationality
criterion in several forms, the definitions and degree bound of the polynomials $T_n$, the minimum
$-1847/1100$ of the prime-$2$ profile for the band construction, and the margin inequality
\texttt{decayRate\_lt\_contentRate} for the rounded rates. The formalisation of the three
estimates is in preparation. According to \path{logs/verification.md}, the package was built with
\texttt{lake exe cache get \&\& lake build} on 25 November 2026 by the AI system that wrote it,
and \texttt{\#print axioms} reported only \texttt{propext}, \texttt{Classical.choice} and
\texttt{Quot.sound} for each of the $29$ declarations of \path{lean/theorems.toml}
(\path{lean/AXIOMS.txt}); the sources contain no \texttt{sorry}, \texttt{admit}, \texttt{axiom}
declaration or \texttt{native\_decide}. We did not rerun the build for this revision of the
manuscript. No independent human review of this work has taken place.

\afbacksection{Acknowledgements}

The construction studied here starts from the work of Shawn Krisman~\cite{Krisman2026chi2,
Krisman2026chi2small}, whose determinants, criterion, reflection argument, size bound and
nonvanishing strategy are reused and adapted throughout; we have tried to indicate each such use
at the place where it occurs.

\afbacksection{AI-assisted research disclosure}

\emph{Exploration and candidate proofs.} According to the owners of AxiomForge, the original
exploration, the choice of the two constructions, the candidate proofs and the first versions
of the certificate programs were produced by GPT-6.2 XHigh (OpenAI), used through its web
interface.

\emph{Reconstruction, audit and manuscript.} Claude Opus 5.5 (Anthropic), working under the
direction of AxiomForge, reconstructed the arguments from the candidate material, rewrote the
proofs, wrote the second implementations of all finite computations (Appendix~\ref{app:certificates}),
carried out the small-size audit (Appendix~\ref{app:consistency}), wrote this manuscript, and wrote the
Lean package described in the section on reproducibility.

\emph{Review and revision.} An adversarial review of an earlier version of this manuscript was
carried out by Claude Fable 5.2 (Anthropic) in a separate session; it is an AI review, neither
independent of the AI systems involved nor a human review, and it is not a certification. The
revision of 26 November 2026 that answers it (literature statements checked against the cited
preprint, corrected statements of \cref{lem:filling} and \cref{prop:profile}, repaired verifier,
reproducibility section) was carried out by Claude Sonnet 5.5 (Anthropic), working under the
direction of AxiomForge.

\emph{Human role and review.} The owners of AxiomForge chose the problem, directed the work and
decided on its publication as a candidate. No human expert independent of the AI systems has
reviewed the mathematics. The checks reported here were performed by AI systems and are not an
independent certification; AI involvement in a step should not be read as evidence of its
correctness.

\emph{Responsibility.} AI systems are not authors of this work and cannot take responsibility
for it. AxiomForge takes full responsibility for the entire content of this manuscript,
including every part produced with AI assistance, and for any errors it may contain. Reports
of errors are welcome at \href{mailto:axiomforge.math@gmail.com}{axiomforge.math@gmail.com};
confirmed errors are corrected openly in a new version, recorded in the changelog of this
contribution, with credit to whoever reported them.


\clearpage
\bibliographystyle{amsplain}
\bibliography{references}

\end{document}
